% !TeX root = ../../Book3.tex
% 纲要标题：## 第10章　长度空间、Gromov–Hausdorff 距离与预紧性
% 纲要位置：几何学计划纲要.md，第 2460 行。
% 正文状态：初稿；尚未正式审核。

% 纲要内容开始
% 
% > 本章目标与证明范围：完整建立度量极限工具，先以可计算退化模型解释收敛，再证明覆盖数控制的预紧性。一般 GH 收敛不蕴含拓扑或光滑结构稳定。
% 
% 纲要内容结束

\chapter{长度空间、Gromov–Hausdorff 距离与预紧性}
\label{g3:ch:10}
\BookChapterNumber{g3:ch:10}

\begin{chapterabstract}
一列几何空间可以在距离意义下趋近，同时发生维数下降或出现分叉。本章从曲线长度与内在距离出发，证明长度空间的 Hopf--Rinow 定理，再构造 Gromov--Hausdorff 距离，并用对应的畸变和近似映射给出可计算的等价描述。通过有限网的距离矩阵、零距离商及完备化，完整证明统一覆盖数的预紧性判据和 GH 等距类空间的完备性。非紧情形采用带基点收敛，计算光滑点切极限、扁平环面的多尺度塌缩和测地树的非唯一切极限；最后用 CAT 与 Alexandrov 比较说明光滑曲率界在度量空间中的表达。
\end{chapterabstract}

\begin{learninggoals}
\begin{itemize}
	\item 区分原距离、内在距离及测地距离，说明完备性与局部紧性如何保证最短路存在。
	\item 通过对应及畸变计算 GH 上、下界，并从零距离构造真正的满射等距。
	\item 用有限网和距离矩阵逐步构造极限，证明统一覆盖数的预紧性判据。
	\item 处理带基点收敛中的球半径误差，计算缩放与塌缩同时发生时的极限。
	\item 检验 CAT 上比较与 Alexandrov 下比较，并说明小三角形如何检测截面曲率。
\end{itemize}
\end{learninggoals}

把周长分别为 \(L\) 和 \(h\) 的两个圆周作平坦乘积，得到面积为 \(Lh\) 的二维环面。固定 \(L\) 而令 \(h\downarrow0\)，两个点之间的距离是两坐标圆周距离的平方和开根。忘掉第二坐标以后，每对点的距离改变至多 \(h/2\)，而第一个圆周的每个点都有原像。于是整个二维环面越来越接近一个一维圆周。这个比较可以直接计算，却无法通过“在同一套二维坐标中度量系数趋于正定矩阵”来描述。

\ProseParagraphBreak

怎样给两个底集合不同、甚至拓扑不同的空间定义距离？若先把它们放进某个公共度量空间，就能比较两者作为子集的 Hausdorff 距离。Gromov--Hausdorff 距离对所有这样的放置取最小可能的误差。等价地，可以直接在两空间的点之间建立对应，要求每对距离的差都小。后一种写法使有限网、缩放、投影和塌缩进入同一种估计。

\ProseParagraphBreak

有界直径并不足以保证一个空间族有收敛子列：越来越多、两两距离为一的点仍可挤进同一个有界直径。真正需要控制的是每个尺度上的覆盖数。证明的主要工作，是把各空间中并不天然对应的网点统一编号，使全部距离矩阵元沿同一个子列收敛，再补上这些标号所逼近的点。非紧空间还要同时记录基点附近的所有整数半径，保证各球极限在一个共同空间中相容。

\ProseParagraphBreak

前四节所需的基本度量拓扑、Cauchy 完备化与紧性判据在正文中恢复。黎曼应用使用第二章的法坐标与完备性，第八章的 Bishop--Gromov 比较提供覆盖数，第九章的空间形式给出平坦商模型。最后一节使用第七章的 Hessian 与 Riccati 公式，以及第九章的三角比较。阅读对照可查 Burago、Burago 与 Ivanov 的第2章“Length Spaces”、第7章“Space of Metric Spaces”、第8章“Large-scale Geometry”，合成曲率的进一步内容见其第9—10章\cite[第2章，第7—10章]{BuragoBuragoIvanov2001}。

% !TeX root = ../../Book3.tex
% 纲要标题：### 10.1 长度空间与测地空间
% 纲要位置：几何学计划纲要.md，第 2464 行。
% 正文状态：初稿；尚未正式审核。

% 纲要内容开始
% 
% * 曲线长度和内在度量
% * 完备、局部紧及测地存在
% * 锥、商和粘合空间
% 
% 纲要内容结束

\section{长度空间与测地空间}
\label{g3:sec:1001}

把一条细圆柱压到它的轴上，圆周方向的距离逐渐消失，极限仍有距离，却不再是原来维数的流形。为了容纳这种变化，先保留曲线长度和最短路，把坐标与切空间暂时放下。第二章的黎曼距离正是由曲线长度定义的；本节把同一思想用于任意度量空间，并说明局部紧性怎样代替有限维坐标中的紧球。

\subsection{曲线长度和内在度量}
\label{g3:subsec:1001:01}
在度量空间中不能先写速度的积分，但可以沿曲线选取有限个点，以相邻点的距离之和逼近长度。这个定义只使用已经存在的距离。

\begin{definition}\label{g3:def:metric-curve-length}
	设 \((X,d)\) 为度量空间，\(\gamma:[a,b]\to X\) 连续。其\term[quxian changdu]{曲线长度}[length of a curve]为
	\[
		L_d(\gamma)=\sup_{a=t_0<\cdots<t_N=b}
		\sum_{i=1}^N d(\gamma(t_{i-1}),\gamma(t_i))\in[0,\infty].
	\]
	长度有限的曲线称为\term[ke chang quxian]{可长曲线}[rectifiable curve]。由 \(d\) 诱导的\term[neizaijuli]{内在距离}[intrinsic distance]定义为
	\[
		d_\ell(x,y)=\inf\{L_d(\gamma):\gamma\text{ 为连接 }x,y\text{ 的连续曲线}\},
	\]
	没有这样的可长曲线时取 \(+\infty\)。若对所有点对都有 \(d=d_\ell\)，则 \((X,d)\) 称为\term[changdukongjian]{长度空间}[length space]。
\end{definition}

\begin{proposition}\label{g3:prop:metric-length-properties}
	设 \((X,d)\) 为度量空间。曲线长度在限制后拼接时相加，在保持方向的连续单调满射重参数化下不变。若同一区间上的曲线 \(\gamma_j\) 逐点收敛到连续曲线 \(\gamma\)，则
	\[
		L_d(\gamma)\le\liminf_jL_d(\gamma_j).
	\]
	每条定义在 \([a,b]\)、长度为 \(L<\infty\) 的连续曲线都可写成 \(\gamma(t)=\widehat\gamma(s(t))\)，其中 \(s(t)=L_d(\gamma|_{[a,t]})\) 连续、非减，\(\widehat\gamma:[0,L]\to X\) 为1-Lipschitz 映射，且每个子区间上的长度等于该区间长度。
\end{proposition}
\begin{proof}
	在任意分割中插入一个指定的参数只会增加折线和。把左右两段的分割合并，再分别逼近两段长度，便得可加性。单调满射把分割对应到原参数区间；删除重复参数后，两个方向都能比较折线和，故长度不变。对固定分割，有限个距离在逐点极限下收敛，因而该分割的折线和不超过 \(\liminf_jL_d(\gamma_j)\)；再对分割取上确界，得到下半连续性。

	由可加性，\(s\) 非减。说明它连续时，固定 \(t\) 和 \(\varepsilon>0\)，选含 \(t\) 的有限分割，使折线和大于 \(L-\varepsilon\)。若 \(u<t\) 足够接近 \(t\)，将 \(u\) 插入这个分割；除 \([u,t]\) 外各段的实际长度不小于相应折线和，于是
	\[
		L_d(\gamma|_{[u,t]})\le\varepsilon+d(\gamma(u),\gamma(t)).
	\]
	令 \(u\uparrow t\)，再令 \(\varepsilon\downarrow0\)，得左连续；右连续同理。因此 \(s\) 连续且映满 \([0,L]\)。若 \(s(t)=s(u)\)，两点间子曲线长为零，所以 \(\gamma(t)=\gamma(u)\)。据此令 \(\widehat\gamma(s(t))=\gamma(t)\)，定义良好，并有
	\[
		d(\widehat\gamma(v),\widehat\gamma(w))\le|v-w|.
	\]
	它的子段长度不超过参数差，原曲线相应子段的长度则恰为该参数差；利用已证的重参数化不变性，得到等号。\(L=0\) 时曲线恒定，单点区间上的定义同样成立。
\end{proof}

拼接曲线说明 \(d_\ell\) 满足三角不等式，而端点距离不超过长度，所以 \(d_\ell\ge d\)。在任意点对可由可长曲线连接的部分，\(d_\ell\) 是有限值度量。它已经是长度度量：对一条原来可长的曲线，任意子段端点的 \(d_\ell\) 距离不超过该子段的原长度，故新旧长度相同；弧长函数的连续性还保证这条曲线对新度量连续。对原来近乎最短的曲线使用这一观察，得到 \((d_\ell)_\ell=d_\ell\)。新度量所给拓扑可能更细，因而这项连续性检验不能省去。

\begin{definition}\label{g3:def:metric-geodesic-space}
	设 \((X,d)\) 为度量空间。若 \(\gamma:[0,1]\to X\) 满足
	\[
		d(\gamma(s),\gamma(t))=|s-t|d(\gamma(0),\gamma(1)),
	\]
	则称它为一条常速\term[cedi duan]{测地线段}[geodesic segment]。任意两点都可用这种线段连接的空间称为\term[cedi kongjian]{测地空间}[geodesic space]。常值线段包括在内。
\end{definition}

本章的测地线段总实现端点距离；黎曼联络所定义的测地线则只在充分短的子段上有此性质。测地空间一定是长度空间，反向结论需要额外条件。

\begin{example}\label{g3:exm:intrinsic-metric-models}
	半径 \(a>0\) 的圆周带环境弦距离时，两个点的距离为 \(2a\sin(\theta/2)\)，其中 \(\theta\in[0,\pi]\) 为较小圆心角。任何绕圆的连接至少走过较短圆弧，故内在距离为 \(a\theta\)。两种距离诱导相同拓扑，但只有后者是长度度量。

	去心平面 \(\mathbb R^2\setminus\{0\}\) 带 Euclid 距离仍是长度空间。若线段穿过原点，把半径 \(\delta\) 内的一段替换成半圆，其额外长度为 \((\pi-2)\delta\)，可任意小。然而 \((-1,0)\) 到 \((1,0)\) 没有长度为2的连接：Euclid 长度等于端点距离时，曲线的全部点须落在线段上，而该线段经过缺失点。因此这个长度空间不是测地空间。
\end{example}

\subsection{完备、局部紧及测地存在}
\label{g3:subsec:1001:02}
为了使越来越短的曲线有收敛子列，既要排除有限距离内的缺失点，也要把它们限制在紧集内。下面先把这些紧性条件写成可直接使用的度量语言。

\begin{definition}\label{g3:def:proper-metric-space}
	设 \((X,d)\) 为度量空间，\(\varepsilon>0\)。空间的每个点若有一个闭包紧的邻域，就称空间局部紧。若每个有限半径闭球都紧，就称它为\term[bi qiu jin]{闭球紧的度量空间}[proper metric space]。对每个 \(\varepsilon>0\)，若空间都能被有限个半径 \(\varepsilon\) 的开球覆盖，就称它\term[quan youjie]{全有界}[totally bounded]。有限集合 \(F\subset X\) 若使每个点到 \(F\) 的距离不超过 \(\varepsilon\)，就称为有限 \(\varepsilon\)-网；本章的网允许边界等号。将半径缩小一半可见，用开球或闭球覆盖检验全有界，条件等价。
\end{definition}

\begin{lemma}\label{g3:lem:complete-totally-bounded}
	度量空间紧，当且仅当它完备且全有界。
\end{lemma}
\begin{proof}
	紧性使每个等半径球覆盖有有限子覆盖，故全有界。紧空间内任意序列都有收敛子列：若没有聚点，每点都有只含有限项的邻域，有限子覆盖便与序列有无限项矛盾；从聚点的逐级小球中可选出收敛子列。Cauchy 序列一旦有收敛子列，整列就收敛，所以空间完备。

	反过来，用有限个半径 \(2^{-m}\) 的球逐级筛选序列：每一级保留一个含有无限项的球，并从中选取递增下标。所得对角子列是 Cauchy 列，故在完备空间中收敛。最后说明序列紧性足以给出开覆盖紧性。序列紧的空间必全有界，否则可在某个固定尺度贪心选出两两分离的点，形成没有收敛子列的序列。若一个开覆盖不存在统一的正数 \(\delta\)，使每个 \(\delta\)-球落在某个覆盖成员中，则对每个 \(m\) 可取 \(x_m\)，使 \(B(x_m,1/m)\) 不包含于任一覆盖成员。取收敛子列及含其极限的开成员，该成员最终包含这些小球，矛盾。因此存在上述 \(\delta>0\)。再取有限 \(\delta/2\)-网，即得到原覆盖的有限子覆盖。
\end{proof}

\begin{proposition}\label{g3:prop:midpoint-geodesic}
	设 \((X,d)\) 完备，且任意 \(x,y\in X\) 都有中点，即存在 \(m\) 满足 \(d(x,m)=d(m,y)=d(x,y)/2\)。则 \(X\) 是测地空间。
\end{proposition}
\begin{proof}
	固定 \(x,y\)，令 \(D=d(x,y)\)。先规定 \(\gamma(0)=x\)、\(\gamma(1)=y\)，再反复在相邻已定点间取中点，从而在所有二进分数上定义 \(\gamma\)。第 \(j\) 级相邻点的距离为 \(D/2^j\)。对同一级的 \(s<t\)，相邻距离相加给出 \(d(\gamma(s),\gamma(t))\le D(t-s)\)；将端点间的总距离 \(D\) 与前、中、后三段比较，又迫使这里取等号。任意两个二进分数可放到共同一级，因此这个等式在稠密参数集上成立。完备性使它唯一延拓到 \([0,1]\)，距离的连续性保持等式，得到所需测地线段。
\end{proof}

\begin{theorem}[长度空间的 Hopf--Rinow 定理]\label{g3:thm:metric-hopf-rinow}
	每个非空、完备、局部紧的长度空间都闭球紧，而且是测地空间。
\end{theorem}
\begin{proof}
	先证明闭球紧。固定 \(p\in X\)，记 \(\mathcal R\) 为使 \(\overline B(p,r)\) 紧的非负半径集合。局部紧性保证其中有正数；紧球的闭子集仍紧，所以 \(\mathcal R\) 向下封闭。

	假设 \(R=\sup\mathcal R<\infty\)。对任意 \(\varepsilon>0\)，取 \(r<R\) 且 \(R-r<\varepsilon/3\)。若 \(x\in\overline B(p,R)\) 而 \(d(p,x)>r\)，取长度小于 \(d(p,x)+\varepsilon/3\) 的连接曲线，按弧长在时间 \(r\) 处截取一点 \(z\)。则 \(z\in\overline B(p,r)\)，并且
	\[
		d(z,x)<R-r+\varepsilon/3<2\varepsilon/3.
	\]
	球 \(\overline B(p,r)\) 紧，有有限 \(\varepsilon/3\)-网；它也是 \(\overline B(p,R)\) 的有限 \(\varepsilon\)-网。因此后者全有界，又是完备空间的闭子集，上一引理使它紧。

	现在把这个紧球再扩大一点。对每个 \(z\in\overline B(p,R)\)，选 \(\eta_z>0\)，使 \(\overline B(z,4\eta_z)\) 紧。有限个 \(B(z_i,\eta_i)\) 覆盖 \(\overline B(p,R)\)，令 \(\eta=\min_i\eta_i\)。对 \(x\in\overline B(p,R+\eta)\)，若 \(d(p,x)>R\)，取长度小于 \(d(p,x)+\eta\) 的连接，在弧长 \(R\) 处取 \(z\in\overline B(p,R)\)。于是 \(d(x,z)<2\eta\)，某个 \(i\) 满足 \(d(z,z_i)<\eta_i\)，从而 \(x\in B(z_i,3\eta_i)\)。原紧球内的点也在这些球中。因此
	\[
		\overline B(p,R+\eta)\subseteq\bigcup_i\overline B(z_i,4\eta_i).
	\]
	右边紧，左边闭，所以左边紧，与 \(R\) 的定义矛盾。故所有有限半径闭球紧。

	固定 \(x,y\)，记 \(D=d(x,y)\)。若 \(D=0\)，取常线。否则选长度 \(L_j<D+1/j\) 的连接曲线，并在弧长一半处取 \(m_j\)。这些点留在固定紧球中，可取子列收敛到 \(m\)，而
	\[
		d(x,m),d(m,y)\le D/2.
	\]
	两项之和不小于 \(D\)，所以二者均为 \(D/2\)。任意点对都有中点，应用\cref{g3:prop:midpoint-geodesic}得到测地存在性。
\end{proof}

这个证明将最短线的存在化成中点的紧性问题。完备性补上二进参数的极限，局部紧性则经长度性质传播到任意有界尺度。若空间本来闭球紧，证明的最后一段就表明：每个闭球紧的长度空间都是测地空间。

\ProseParagraphBreak

各条件有不同作用。去心平面的例子缺少完备性。无限维 Hilbert 空间虽完备且任意两点间都有直线，闭单位球却不紧：一列正交单位向量的两两距离为 \(\sqrt2\)，无法取收敛子列。它也不局部紧。相反，任意无限集合赋予不同点距离都为1的度量，完备且局部紧，但整个空间是一个非紧的有界闭球；这里缺少长度空间条件。

\subsection{锥、商和粘合空间}
\label{g3:subsec:1001:03}
流形退化后可能出现锥顶、边界或分叉点。下面从距离本身构造这些模型；构造是否给出度量、是否有最短路，都可以直接检验。

\begin{example}\label{g3:exm:flat-metric-cone}
	取 \(\theta>0\)，把 \([0,\infty)\times(\mathbb R/\theta\mathbb Z)\) 的全部零半径点识别为一个点 \(o\)。在 \(r>0\) 上按
	\(\mathrm dr^2+r^2\mathrm d\varphi^2\)
	计算曲线长度，并允许曲线经过 \(o\)。取这些长度的下确界得到\term[pingtan yuanzhui]{平坦圆锥}[flat cone] \(C_\theta\)。若 \(\delta\in[0,\theta/2]\) 是两个角坐标的圆周距离，则
	\begin{equation}\label{g3:eq:flat-cone-distance}
		d_{C_\theta}((r,\varphi),(s,\psi))^2
		=r^2+s^2-2rs\cos\bigl(\min\{\delta,\pi\}\bigr).
	\end{equation}
	当 \(\delta\ge\pi\) 时，右侧为 \((r+s)^2\)。

	说明公式时，先看不经过锥顶的曲线。角坐标可以沿参数区间连续提升到实数，角度总变化的绝对值至少为 \(\delta\)。若该变化小于 \(\pi\)，平面展开后的曲线长不小于两端点弦长；在所扫过的角扇形内用直线段实现这个下界。若角度变化的绝对值至少为 \(\pi\)，在首次转过 \(\pi\) 处记半径为 \(u\)。前段展开后端点位于相反射线上，长度至少为 \(r+u\)；后段长度至少为 \(|s-u|\)，总长至少为 \(r+s\)。经过锥顶的曲线也至少长 \(r+s\)，而两条径向线拼接达到这一值。对允许的角度提升取最小，便得到公式；\(\delta<\pi\) 时若采用更大的提升角，所得下界不会更小。

	长度拼接保证三角不等式，公式保证除已识别的零半径点外距离为正。每个锥顶闭球是紧集 \([0,R]\times(\mathbb R/\theta\mathbb Z)\) 的连续像，其他有界闭球为更大锥顶紧球的闭子集。因此 \(C_\theta\) 闭球紧且测地。伸缩 \((r,\varphi)\mapsto(\lambda r,\varphi)\) 把所有距离乘以 \(\lambda\)；当 \(\theta=2\pi\) 时得到 Euclid 平面。角度改变后，锥顶附近小圆周的长度是 \(\theta r\)，显示了与平面的差别。
\end{example}

\begin{proposition}\label{g3:prop:compact-group-metric-quotient}
	设紧拓扑群 \(G\) 连续地以等距变换作用于度量空间 \(X\)。在轨道集合上定义
	\[
		d_{X/G}([x],[y])=\min_{g\in G}d_X(x,gy).
	\]
	则这是一个度量，商投影为1-Lipschitz 映射，且诱导商拓扑。若 \(X\) 是长度空间或测地空间，商也分别具有这一性质。若 \(X\) 闭球紧，则商闭球紧。
\end{proposition}
\begin{proof}
	紧性和连续性使最小值取得；更换轨道代表不改变最小值。若距离为零，则两点在同一轨道。对任意 \(g,h\in G\)，
	\[
		d(x,ghz)\le d(x,gy)+d(gy,ghz)
		=d(x,gy)+d(y,hz),
	\]
	取最小值即得三角不等式。取恒等元得到投影不增加距离。又有
	\(\pi(B(x,r))=B([x],r)\)，
	所以投影为开映射，度量拓扑恰为商拓扑。

	取达到轨道距离的代表 \(gy\)，把 \(x\) 到 \(gy\) 的近乎最短曲线投到商上。投影不增加长度，曲线长度又不小于商中端点距离，故商是长度空间。若原空间测地，可投影一条长恰为商距离的最短段；其所有子段也实现距离，否则替换该子段就会使总长变短。按常速参数化便得到商测地线。最后，最小值取得还说明
	\[
		\overline B([x],R)=\pi(\overline B(x,R)).
	\]
	右边在原空间闭球紧时为紧集，完成证明。
\end{proof}

例如，平面按旋转角 \(2\pi/m\) 的有限循环群取商，便得到角度 \(2\pi/m\) 的圆锥；按关于一条直线的反射取商，得到闭半平面。稳定点可以产生锥顶或边界，因而商空间不必仍为无边界流形。

\begin{proposition}\label{g3:prop:metric-gluing-along-closed-set}
	设 \(X,Y\) 为度量空间，\(A\subset X\)、\(B\subset Y\) 非空闭，\(f:A\to B\) 为满射等距。把 \(a\) 与 \(f(a)\) 识别所得集合 \(Z=X\cup_fY\) 上，保留每一侧原距离，并对跨侧点规定
	\begin{equation}\label{g3:eq:metric-gluing-distance}
		d_Z(x,y)=\inf_{a\in A}\{d_X(x,a)+d_Y(f(a),y)\},
		\qquad x\in X,\ y\in Y.
	\end{equation}
	这些公式给出度量，使两侧等距嵌入。若两侧都是长度空间，则 \(Z\) 是长度空间；若两侧闭球紧，则 \(Z\) 闭球紧。因此两侧闭球紧且为长度空间时，粘合空间也测地。
\end{proposition}
\begin{proof}
	当一个点属于粘合子集时，跨侧公式与同侧距离相符，因为 \(f\) 等距且可直接取该点为中转点。检查三角不等式的唯一新情况，是从一侧出发经另一侧返回。例如对 \(x,x'\in X\)、\(y\in Y\)，任取 \(a,b\in A\)，有
	\[
	\begin{aligned}
		d_X(x,x')&\le d_X(x,a)+d_X(a,b)+d_X(b,x')\\
		&=d_X(x,a)+d_Y(f(a),f(b))+d_X(b,x')\\
		&\le d_X(x,a)+d_Y(f(a),y)+d_Y(y,f(b))+d_X(b,x').
	\end{aligned}
	\]
	对 \(a,b\) 取下确界即得所需不等式，其余跨侧情形用同一侧的三角不等式即可。若跨侧距离为零，可取 \(a_j\) 使两项之和趋零；于是 \(a_j\to x\)，闭性给出 \(x\in A\)，再由 \(f(a_j)\to f(x)\) 得 \(y=f(x)\)。所以零距离恰对应已经作出的识别。

	两侧原距离未改变。两侧为长度空间时，同侧近乎最短曲线直接保留；跨侧先选接近下确界的 \(a\)，再拼接两侧近乎最短曲线，便得长度性质。

	最后固定 \(a_0\in A\)。对 \(x\in X\)，有 \(d_Z(a_0,x)=d_X(a_0,x)\)，对 \(y\in Y\) 同理。因此以粘合点为中心的闭球，恰为两侧相应闭球的并。在两侧闭球紧时，这个有限并紧；任意中心的有界闭球包含在某个这样的紧球中且为闭集，所以也紧。测地结论由\cref{g3:thm:metric-hopf-rinow}的闭球紧版本得到。
\end{proof}

将若干区间的一个端点粘在一起会产生分叉，而两张平面沿直线等距粘合会产生一条公共接缝。若粘合子集不闭，零距离还可能识别原来没有声明相等的点：把两份实直线沿开半轴 \((0,\infty)\) 粘合，两份原点在式\eqref{g3:eq:metric-gluing-distance}下的距离为零。这时应再取零距离商，不能直接把原集合上的公式称为度量。
% 纲要条目排为编号小节标题，并进入目录与书签；说明段仍保留于纲要注释中。

% !TeX root = ../../Book3.tex
% 纲要标题：### 10.2 GH 距离的构造
% 纲要位置：几何学计划纲要.md，第 2470 行。
% 正文状态：初稿；尚未正式审核。

% 纲要内容开始
% 
% * 公共度量空间中的 Hausdorff 距离
% * 对应与畸变
% * ε-逼近映射的等价描述
% * 零距离与等距类
% 
% 纲要内容结束

\section{GH 距离的构造}
\label{g3:sec:1002}

同一个圆可以平移到平面中很远的地方，圆本身的距离结构却没有改变。另一方面，一条越来越细的圆柱即使没有与区间相同的点集，沿轴投影后仍能逐渐保留全部距离。这提示我们：比较空间时，应允许重新安放两个空间，并用一个点对应几个相近的点。本节把这两种想法分别写成公共空间中的距离和对应关系中的畸变，再证明它们给出同一个量。

\subsection{公共度量空间中的 Hausdorff 距离}
\label{g3:subsec:1002:01}

先在固定的度量空间 \((Z,d_Z)\) 中比较两个非空紧子集。点到集合的距离记为
\[
	d_Z(z,A)=\inf_{a\in A}d_Z(z,a).
\]
由三角不等式，\(z\mapsto d_Z(z,A)\) 是1-Lipschitz 函数。当 \(A\) 紧时，下确界在某个点取得；这样的最近点可以不唯一。若只取 \(\inf_{a\in A,b\in B}d_Z(a,b)\)，两个相交但相差很大的集合也会得到零。为了测量整个集合是否接近，需要要求每个点都能在另一边找到近点。

\begin{definition}\label{g3:def:hausdorff-distance}
	设 \(A,B\) 是度量空间 \((Z,d_Z)\) 的非空紧子集。定义它们的\term[Hausdorff juli]{Hausdorff 距离}[Hausdorff distance]为
	\begin{equation}\label{g3:eq:hausdorff-distance}
		d_H^Z(A,B)=\max\left\{\sup_{a\in A}d_Z(a,B),\ \sup_{b\in B}d_Z(b,A)\right\}.
	\end{equation}
	等价地，它是满足 \(A\subset U_r(B)\)、\(B\subset U_r(A)\) 的 \(r>0\) 的下确界，其中 \(U_r(A)=\{z:d_Z(z,A)<r\}\)。
\end{definition}

两个有向距离都要出现：\(A\subset B\) 只能使第一个上确界为零，第二个还要测量 \(B\) 中多出的部分。这个距离中的 Hausdorff 与拓扑分离条件中的人物相同。\footnote{Felix Hausdorff（1868--1942），德国数学家，研究集合论、拓扑学与测度论；其工作奠定了拓扑空间、度量空间及分数维数理论的重要基础。}

\begin{proposition}\label{g3:prop:hausdorff-metric}
	设 \((Z,d_Z)\) 为度量空间。\(d_H^Z\) 在 \(Z\) 的非空紧子集全体上是有限值度量。
\end{proposition}
\begin{proof}
	紧集有界，所以式\eqref{g3:eq:hausdorff-distance}有限。非负性和对称性直接成立。若 \(d_H^Z(A,B)=0\)，则每个 \(a\in A\) 都属于 \(\overline B=B\)，故 \(A\subset B\)；反向同理，于是 \(A=B\)。这里用到紧子集在度量空间中闭。

	设 \(A,B,C\) 都非空且紧。对每个 \(a\in A\)，取 \(b\in B\) 实现 \(d_Z(a,B)\)，再取 \(c\in C\) 实现 \(d_Z(b,C)\)，则
	\[
		d_Z(a,C)\le d_Z(a,b)+d_Z(b,c)
		\le d_H^Z(A,B)+d_H^Z(B,C).
	\]
	对 \(a\) 取上确界，并交换 \(A,C\) 作同样估计，得到三角不等式。
\end{proof}

\begin{example}\label{g3:exm:hausdorff-position}
	在实直线中，\(A=[0,1]\)、\(B=[L,L+1]\)，其中 \(L\ge0\)，满足 \(d_H^{\mathbb R}(A,B)=L\)。平移 \(x\mapsto x+L\) 又是 \(A\) 到 \(B\) 的等距映射。因而 Hausdorff 距离保留了集合在环境中的位置；要只比较两个空间自身的度量，还须允许改变这两个等距副本的位置。
\end{example}

\begin{definition}\label{g3:def:gh-distance}
	设 \((X,d_X)\)、\((Y,d_Y)\) 是非空紧度量空间。它们的\term[Gromov Hausdorff juli]{Gromov--Hausdorff 距离}[Gromov--Hausdorff distance]定义为
	\begin{equation}\label{g3:eq:gh-distance-ambient}
		d_{\mathrm{GH}}(X,Y)
		=\inf_{Z,\iota_X,\iota_Y}
		d_H^Z\bigl(\iota_X(X),\iota_Y(Y)\bigr),
	\end{equation}
	其中 \(Z\) 遍历度量空间，\(\iota_X:X\to Z\)、\(\iota_Y:Y\to Z\) 遍历保持距离的嵌入。简称为 GH 距离。
\end{definition}
\symidx{\(d_{\mathrm{GH}}(X,Y)\)：非空紧度量空间之间的 Gromov--Hausdorff 距离}

保持距离要求 \(d_Z(\iota_Xx,\iota_Xx')=d_X(x,x')\)，因此不能为了使两个副本接近而改变各自内部的距离。公共空间也不必是 Euclid 空间；这一自由度正是下面的构造能够工作的原因。

\ProseParagraphBreak

定义中的候选总存在。取不交并 \(X\sqcup Y\)，保留两边原距离，并令所有跨边距离等于 \(D=\max\{1,\operatorname{diam}X,\operatorname{diam}Y\}\)。跨边三角不等式由两边直径不超过 \(D\) 保证，从而得到一个度量。因此 GH 距离有限。若两个空间等距，可将它们嵌入同一个副本，故距离为零。反向结论将在本节末证明。

\subsection{对应与畸变}
\label{g3:subsec:1002:02}

公共环境的定义有直接的几何含义，但寻找所有环境不便计算。若两个副本在环境中很近，可以把彼此相近的点配对；反过来，若配对保持距离到一个小误差，也可以在每对点之间架设短桥。这样便把环境的选择转化为一个仅涉及两边距离表的问题。

\begin{definition}\label{g3:def:metric-correspondence}
	设 \(X,Y\) 是非空集合。若子集 \(R\subset X\times Y\) 向两个因子的投影都满射，即每个 \(x\in X\) 至少与一个 \(y\in Y\) 配对，每个 \(y\in Y\) 也至少与一个 \(x\in X\) 配对，则称 \(R\) 为 \(X,Y\) 之间的一个\term[duiying]{对应}[correspondence]。
	若 \(X,Y\) 还带有度量，则定义对应的\term[jibian]{畸变}[distortion]为
	\begin{equation}\label{g3:eq:correspondence-distortion}
		\operatorname{dis}R
		=\sup_{(x,y),(x',y')\in R}
		\bigl|d_X(x,x')-d_Y(y,y')\bigr|.
	\end{equation}
\end{definition}

对应允许一对多或多对一。一条满射的图像是对应，任意映射的图像则可能漏掉靶空间中的点。若 \((x,y),(x',y)\in R\)，畸变条件给出 \(d_X(x,x')\le\operatorname{dis}R\)：被配到同一点的整簇点必须足够接近。这正好适合描述塌缩。

\begin{theorem}\label{g3:thm:gh-correspondence}
	对任意非空紧度量空间 \(X,Y\)，有
	\begin{equation}\label{g3:eq:gh-correspondence}
		d_{\mathrm{GH}}(X,Y)=\frac12\inf_R\operatorname{dis}R,
	\end{equation}
	其中下确界取遍 \(X,Y\) 之间的全部对应。
\end{theorem}
\begin{proof}
	先从公共空间构造对应。设 \(X,Y\) 已等距嵌入 \(Z\)，并取 \(r>d_H^Z(X,Y)\)。集合
	\[
		R=\{(x,y):d_Z(x,y)<r\}
	\]
	向两边都满射。对其中任意两对点，由环境中的三角不等式得
	\[
		\bigl|d_X(x,x')-d_Y(y,y')\bigr|
		\le d_Z(x,y)+d_Z(x',y')<2r.
	\]
	故 \(\operatorname{dis}R\le2r\)。先令 \(r\) 降至该 Hausdorff 距离，再对环境取下确界，得到式\eqref{g3:eq:gh-correspondence}右端不大于左端。

	反过来，给定对应 \(R\)，记 \(\Delta=\operatorname{dis}R\)，并取 \(\eta>\Delta/2\)。在 \(X\sqcup Y\) 上保留两边的原距离，对 \(x\in X,y\in Y\) 定义
	\begin{equation}\label{g3:eq:gh-bridge-metric}
		D(x,y)=D(y,x)
		=\inf_{(a,b)\in R}
		\bigl(d_X(x,a)+\eta+d_Y(b,y)\bigr).
	\end{equation}
	可以把括号内的三项理解为先在 \(X\) 中走到 \(a\)，经过长为 \(\eta\) 的桥，再在 \(Y\) 中走到 \(y\)。\(R\ne\varnothing\) 保证此值有限，而且跨边距离至少为 \(\eta>0\)。

	现在核对三角不等式，尤其要排除经过另一边的两座桥反而缩短原来的距离。若两个端点在 \(X\)，中间点 \(y\) 在 \(Y\)，则对任意 \((a,b),(a',b')\in R\)，有
	\[
	\begin{aligned}
		&d_X(x,a)+\eta+d_Y(b,y)
		+d_Y(y,b')+\eta+d_X(a',x')\\
		&\quad\ge d_X(x,a)+d_Y(b,b')+2\eta+d_X(a',x')\\
		&\quad\ge d_X(x,a)+d_X(a,a')+(2\eta-\Delta)+d_X(a',x')\\
		&\quad\ge d_X(x,x').
	\end{aligned}
	\]
	分别取下确界，即得 \(D(x,y)+D(y,x')\ge d_X(x,x')\)。交换 \(X,Y\) 得到另一种同类不等式。余下的混合情形形如
	\[
		D(x,y)\le d_X(x,x')+D(x',y),
	\]
	它由 \(d_X(x,a)\le d_X(x,x')+d_X(x',a)\) 逐项代入式\eqref{g3:eq:gh-bridge-metric}后取下确界得到。三点都在同一边时使用原距离的三角不等式。对称性、正定性也已成立，所以 \(D\) 确为度量，两边的包含映射等距。

	对每个 \((a,b)\in R\)，直接取这对点作为桥，得 \(D(a,b)\le\eta\)。对应向两边满射，因而 \(d_H^{(X\sqcup Y,D)}(X,Y)\le\eta\)。令 \(\eta\downarrow\Delta/2\)，再对 \(R\) 取下确界，得到左端不大于右端。
\end{proof}

桥长取严格正数，即使 \(\Delta=0\) 也先保留一个正的余量。这样得到的直接就是度量，无须在不交并中再把零距离点作商。公式中的 \(1/2\) 来自两个配对点各自跨一次桥：一个距离误差最多由两段桥长贡献。

\begin{proposition}\label{g3:prop:correspondence-composition}
	设 \(X,Y,Z\) 是非空紧度量空间，\(R\subset X\times Y\)、\(S\subset Y\times Z\) 为对应。定义
	\[
		S\circ R=\{(x,z):\text{存在 }y\in Y,\ (x,y)\in R,\ (y,z)\in S\}.
	\]
	则 \(S\circ R\) 为 \(X,Z\) 之间的对应，并且
	\[
		\operatorname{dis}(S\circ R)\le\operatorname{dis}R+\operatorname{dis}S.
	\]
	从而 GH 距离满足三角不等式。
\end{proposition}
\begin{proof}
	对给定 \(x\)，先用 \(R\) 找到 \(y\)，再用 \(S\) 找到 \(z\)；从给定 \(z\) 反向选择也成立，所以两个投影满射。对 \((x,z),(x',z')\in S\circ R\)，分别选取中间点 \(y,y'\)，则
	\[
	\begin{aligned}
		|d_X(x,x')-d_Z(z,z')|
		&\le |d_X(x,x')-d_Y(y,y')|\\
		&\quad+|d_Y(y,y')-d_Z(z,z')|\\
		&\le\operatorname{dis}R+\operatorname{dis}S.
	\end{aligned}
	\]
	取上确界便得畸变界。对 \(R,S\) 分别取任意接近其畸变下确界的对应，再用\cref{g3:thm:gh-correspondence}，得到
	\(d_{\mathrm{GH}}(X,Z)\le d_{\mathrm{GH}}(X,Y)+d_{\mathrm{GH}}(Y,Z)\)。
\end{proof}

\begin{proposition}\label{g3:prop:gh-diameter}
	设 \(X,Y\) 是非空紧度量空间。则
	\begin{equation}\label{g3:eq:gh-diameter-bound}
		|\operatorname{diam}X-\operatorname{diam}Y|
		\le2d_{\mathrm{GH}}(X,Y).
	\end{equation}
	若 \(\{*\}\) 是单点空间，则
	\[
		d_{\mathrm{GH}}(X,\{*\})=\tfrac12\operatorname{diam}X.
	\]
\end{proposition}
\begin{proof}
	给定对应 \(R\)，对任意 \(x,x'\in X\)，选 \((x,y),(x',y')\in R\)，于是
	\(d_X(x,x')\le\operatorname{diam}Y+\operatorname{dis}R\)。
	取上确界，再交换两边，得到直径之差的绝对值不超过 \(\operatorname{dis}R\)。对应公式给出式\eqref{g3:eq:gh-diameter-bound}。与单点空间之间只有一个对应 \(X\times\{*\}\)，其畸变恰为 \(\operatorname{diam}X\)，所以第二个等式也由对应公式得到。
\end{proof}

\begin{example}\label{g3:exm:gh-scaling-interval}
	设 \(X\) 是非空紧度量空间，\(\lambda>0\)，以 \(\lambda X\) 表示同一点集配距离 \(\lambda d_X\)。恒等配对的畸变为 \(|\lambda-1|\operatorname{diam}X\)，而直径界给出相反方向的同一个数，故
	\[
		d_{\mathrm{GH}}(X,\lambda X)
		=\tfrac12|\lambda-1|\operatorname{diam}X.
	\]
	特别地，对通常距离下的闭区间，\(a,b\ge0\) 时
	\[
		d_{\mathrm{GH}}([0,a],[0,b])=\tfrac12|a-b|.
	\]
	当 \(a,b>0\) 时用比例配对；某个长度为零时用单点公式。也可以把两个区间放在实直线上并使中点重合，直接看到这个 Hausdorff 上界。
\end{example}

\subsection{\texorpdfstring{\(\varepsilon\)-逼近映射的等价描述}{\GeometryPdfEpsilon{}-逼近映射的等价描述}}
\label{g3:subsec:1002:03}

具体退化模型往往已经带有一个自然映射，例如把细圆柱投影到轴线。此时检查每对点的距离误差及像是否覆盖整个目标，比直接设计一个多值对应更方便。这里需要的是加性误差；它允许比误差尺度更小的部分缩成一点。

\begin{definition}\label{g3:def:epsilon-approximation}
	设 \(X,Y\) 是非空紧度量空间，\(\varepsilon>0\)。映射 \(f:X\to Y\) 的畸变定义为
	\[
		\operatorname{dis}f
		=\sup_{x,x'\in X}|d_Y(f(x),f(x'))-d_X(x,x')|.
	\]
	若 \(\operatorname{dis}f\le\varepsilon\)，且
	\[
		\sup_{y\in Y}d_Y(y,f(X))\le\varepsilon,
	\]
	则称 \(f\) 为一个 \(\varepsilon\)-\term[bijin yingshe]{逼近映射}[approximation map]。后一条件称为像在 \(Y\) 中 \(\varepsilon\)-稠密。这里不要求 \(f\) 连续。
\end{definition}

距离畸变为零只保证映射等距嵌入；像稠密到误差尺度的条件保证没有漏掉靶空间的大块区域。\(f(X)\) 不一定闭，故点到像的距离仍按下确界理解。

\begin{proposition}\label{g3:prop:gh-approximation}
	设 \(X,Y\) 是非空紧度量空间，\(\varepsilon>0\)。则：
	\begin{enumerate}
		\item 若 \(d_{\mathrm{GH}}(X,Y)<\varepsilon\)，则存在 \(2\varepsilon\)-逼近映射 \(f:X\to Y\)。交换 \(X,Y\) 后同样成立。
		\item 若存在 \(\varepsilon\)-逼近映射 \(f:X\to Y\)，则 \(d_{\mathrm{GH}}(X,Y)\le3\varepsilon/2\)。
	\end{enumerate}
	因此，对非空紧度量空间序列 \((X_j)\) 及非空紧度量空间 \(X\)，\(d_{\mathrm{GH}}(X_j,X)\to0\) 当且仅当存在 \(\varepsilon_j\to0\) 及 \(\varepsilon_j\)-逼近映射 \(X_j\to X\)。
\end{proposition}
\begin{proof}
	第一项中，由对应公式取 \(R\subset X\times Y\)，使 \(\operatorname{dis}R<2\varepsilon\)。对每个 \(x\) 选一个 \(f(x)\) 使 \((x,f(x))\in R\)。其图像包含于 \(R\)，所以 \(\operatorname{dis}f<2\varepsilon\)。对任意 \(y\in Y\)，选 \(x\) 使 \((x,y)\in R\)，再比较 \((x,y)\) 与 \((x,f(x))\)，得
	\[
		d_Y(y,f(x))=|d_Y(y,f(x))-d_X(x,x)|<2\varepsilon.
	\]
	这证明像的稠密性。

	第二项中，固定任意 \(\delta>0\)，定义
	\[
		R_\delta=\{(x,y):d_Y(f(x),y)\le\varepsilon+\delta\}.
	\]
	它包含 \(f\) 的图像，故向 \(X\) 满射；像的 \(\varepsilon\)-稠密性和余量 \(\delta\) 保证向 \(Y\) 也满射。若 \((x,y),(x',y')\in R_\delta\)，则
	\[
	\begin{aligned}
		|d_X(x,x')-d_Y(y,y')|
		&\le |d_X(x,x')-d_Y(f(x),f(x'))|\\
		&\quad+d_Y(f(x),y)+d_Y(f(x'),y')\\
		&\le3\varepsilon+2\delta.
	\end{aligned}
	\]
	所以 \(d_{\mathrm{GH}}(X,Y)\le3\varepsilon/2+\delta\)。令 \(\delta\downarrow0\)，得到所述常数。最后一项分别应用这两条估计；必要时在 \(d_{\mathrm{GH}}(X_j,X)\) 上增加 \(1/j\) 的严格余量。
\end{proof}

\begin{example}\label{g3:exm:gh-finite-sampling}
	取 \(X=[0,1]\)，\(F_m=\{0,1/m,\ldots,1\}\)，都使用实直线的限制距离。每个点到最近网格点的距离至多 \(1/(2m)\)，故
	\[
		d_{\mathrm{GH}}(F_m,[0,1])\le\frac1{2m}.
	\]
	最近网格点选择 \(q_m:[0,1]\to F_m\) 在相邻网格点的中点发生跳跃，但
	\[
		\bigl||q_m(x)-q_m(x')|-|x-x'|\bigr|
		\le|q_m(x)-x|+|q_m(x')-x'|\le\frac1m.
	\]
	它还满射，因而是 \(1/m\)-逼近。任何从连通区间到有限离散空间的连续映射都是常值；允许这种小幅跳跃，才能用有限样本逼近一个连续空间。
\end{example}

\subsection{零距离与等距类}
\label{g3:subsec:1002:04}

小误差映射可以不连续，逐点选择也不相容。距离为零时，我们要从一列误差趋零的映射中恢复一个真正保持所有距离的映射。紧性使可数稠密集上的像能够同时取极限，再用距离控制填满其余点。

\begin{theorem}\label{g3:thm:gh-isometry-metric}
	对非空紧度量空间 \(X,Y\)，\(d_{\mathrm{GH}}(X,Y)=0\) 当且仅当 \(X,Y\) 等距。因此 GH 距离在非空紧度量空间的等距类上给出度量。
\end{theorem}
\begin{proof}
	等距时距离为零已经由定义证明；有限性、非负性、对称性与三角不等式也已建立。只须证明零距离产生满射等距。

	设 \(d_{\mathrm{GH}}(X,Y)=0\)。由\cref{g3:prop:gh-approximation}取 \(\varepsilon_j\)-逼近 \(f_j:X\to Y\)，其中 \(\varepsilon_j\to0\)。对每个正整数 \(m\)，紧性使 \(X\) 中半径 \(1/m\) 的开球覆盖有有限子覆盖；收集这些有限个中心，得到一个可数稠密子集 \(A=\{a_1,a_2,\ldots\}\)，有限情形允许重复枚举。

	由于 \(Y\) 紧，先取子列使 \(f_j(a_1)\) 收敛，再从中取子列使 \(f_j(a_2)\) 收敛，依次进行。取第 \(j\) 个子列的第 \(j\) 项得到对角子列，仍记为 \(f_j\)，使每个 \(a_i\) 的像都收敛。记其极限为 \(F(a_i)\)。对任意 \(a_i,a_l\)，由畸变趋零得
	\[
		d_Y(F(a_i),F(a_l))
		=\lim_jd_Y(f_j(a_i),f_j(a_l))=d_X(a_i,a_l).
	\]
	所以 \(F\) 在稠密子集上保持距离。

	给定 \(x\in X\)，取 \(a_{i_m}\to x\)。其像是 \(Y\) 中的 Cauchy 序列，而紧度量空间完备，故像有极限，定义为 \(F(x)\)。若改取另一逼近序列，则两列对应项的距离等于其在 \(X\) 中的距离并趋于零，所以极限与选择无关。对两个点分别取逼近序列，距离等式在极限中保持，故延拓后的 \(F:X\to Y\) 等距。

	还需证明满射。事实上 \(f_j\) 沿所选子列一致趋于 \(F\)。固定 \(\delta>0\)，从 \(A\) 中取有限个点 \(a_{i_1},\ldots,a_{i_N}\)，其 \(\delta\)-开球覆盖 \(X\)。对任意 \(x\)，选一个满足 \(d_X(x,a_{i_l})<\delta\) 的中心，则
	\[
	\begin{aligned}
		d_Y(f_j(x),F(x))
		&\le d_Y(f_j(x),f_j(a_{i_l}))
		+d_Y(f_j(a_{i_l}),F(a_{i_l}))
		+d_Y(F(a_{i_l}),F(x))\\
		&\le2\delta+\varepsilon_j+
		\max_{1\le l\le N}d_Y(f_j(a_{i_l}),F(a_{i_l})).
	\end{aligned}
	\]
	右侧的最后两项趋于零。先取 \(j\to\infty\)，再令 \(\delta\downarrow0\)，即得一致收敛。

	对任意 \(y\in Y\)，由像的稠密性选 \(x_j\in X\)，使 \(d_Y(y,f_j(x_j))\le\varepsilon_j+1/j\)。于是
	\[
		d_Y(y,F(X))\le\varepsilon_j+1/j+
		\sup_{x\in X}d_Y(f_j(x),F(x))\longrightarrow0.
	\]
	等距映射连续，所以 \(F(X)\) 紧且闭；这迫使 \(y\in F(X)\)。故 \(F\) 为满射等距。
\end{proof}

\begin{corollary}\label{g3:cor:gh-limit-unique}
	若非空紧度量空间序列 \((X_j)\) 同时 GH 收敛到非空紧度量空间 \(X,Y\)，则 \(X,Y\) 等距。
\end{corollary}
\begin{proof}
	三角不等式给出
	\(d_{\mathrm{GH}}(X,Y)\le d_{\mathrm{GH}}(X,X_j)+d_{\mathrm{GH}}(X_j,Y)\to0\)，再应用上一条定理。
\end{proof}

此后“GH 极限唯一”均指等距意义下唯一。距离趋零同时给出任意有限点组的距离误差趋零，以及整个空间在误差尺度上的覆盖；这两方面将在下一节分别表现为距离矩阵的收敛与有限网的控制。
% 纲要条目排为编号小节标题，并进入目录与书签；说明段仍保留于纲要注释中。

% !TeX root = ../../Book3.tex
% 纲要标题：### 10.3 预紧性和极限构造
% 纲要位置：几何学计划纲要.md，第 2477 行。
% 正文状态：初稿；尚未正式审核。

% 纲要内容开始
% 
% * 有限网及距离矩阵
% * 对角子列构造极限
% * 统一覆盖数的预紧性判据
% * GH 距离的完备性与测地性的保持
% 
% 纲要内容结束

\section{预紧性和极限构造}
\label{g3:sec:1003}

一个紧空间可以由有限个采样点及它们之间的距离近似。若要让一族空间同时具有收敛子列，关键在于同一精度下能否统一控制所需采样点数。距离矩阵把几何对象转成有限个有界实数；逐层加密采样并抽取对角子列，就能在没有预先给定环境空间的情况下构造极限。

\subsection{有限网及距离矩阵}
\label{g3:subsec:1003:01}

先把“有限个点足以近似整个空间”写成定量语言。以下的网采用闭球约定，因而距离恰等于指定半径的点仍被覆盖。

\begin{definition}\label{g3:def:metric-covering-number}
	设 \((X,d)\) 是非空度量空间，\(\varepsilon>0\)。若有限集 \(F\subset X\) 对所有 \(x\in X\) 都满足 \(d(x,F)\le\varepsilon\)，就称 \(F\) 为 \(X\) 的有限 \(\varepsilon\)-网。记
	\[
		N_X(\varepsilon)=\min\{\#F:F\subset X\text{ 是有限 }\varepsilon\text{-网}\},
	\]
	若不存在这样的有限集，则规定 \(N_X(\varepsilon)=\infty\)。这个数称为 \(X\) 在尺度 \(\varepsilon\) 的\term[fugai shu]{覆盖数}[covering number]。因此，\(X\) 全有界等价于对每个 \(\varepsilon>0\) 都有 \(N_X(\varepsilon)<\infty\)。
\end{definition}
\symidx{\(N_X(\varepsilon)\)：用中心在 \(X\) 内的闭 \(\varepsilon\)-球覆盖 \(X\) 所需的最少球数}

紧度量空间的每个尺度都有有限网：用半径 \(\varepsilon\) 的开球覆盖，再抽取有限子覆盖即可。全有界控制精度提高时仍只需要有限个点；它与完备性合在一起，才由\cref{g3:lem:complete-totally-bounded}保证紧性。

\begin{proposition}\label{g3:prop:finite-net-gh-bound}
	设 \(X\) 为非空紧度量空间，\(\varepsilon>0\)，\(F\subset X\) 是有限 \(\varepsilon\)-网，并赋予限制距离。则
	\[
		d_{\mathrm{GH}}(X,F)\le\varepsilon.
	\]
	其中 \(N_X(r)\) 是覆盖数，即半径 \(r\) 的有限网所需最小点数。若非空紧度量空间 \(X,Y\) 满足 \(d_{\mathrm{GH}}(X,Y)<\delta\)，则对每个 \(r>0\)，有
	\begin{equation}\label{g3:eq:covering-number-transfer}
		N_Y(r+2\delta)\le N_X(r).
	\end{equation}
\end{proposition}
\begin{proof}
	第一项直接在公共空间 \(X\) 中计算 Hausdorff 距离：\(F\subset X\)，且每个 \(x\in X\) 到 \(F\) 的距离不超过 \(\varepsilon\)。

	对第二项，由\cref{g3:thm:gh-correspondence}选取畸变小于 \(2\delta\) 的对应 \(\mathcal R\subset X\times Y\)。取一个含 \(N_X(r)\) 个点的 \(r\)-网 \(\{x_1,\ldots,x_m\}\)，并为每个 \(x_j\) 选择 \(y_j\) 使 \((x_j,y_j)\in\mathcal R\)。任给 \(y\in Y\)，先选 \(x\in X\) 使 \((x,y)\in\mathcal R\)，再选 \(j\) 使 \(d_X(x,x_j)\le r\)。畸变控制给出
	\[
		d_Y(y,y_j)<d_X(x,x_j)+2\delta\le r+2\delta.
	\]
	故这些 \(y_j\) 形成 \((r+2\delta)\)-网。重复的 \(y_j\) 只会减少点数，得到所述不等式。
\end{proof}

点数得到控制以后，有限空间的自由度只剩距离矩阵。极限矩阵中两个不同标号之间的距离可能变为零，所以构造时要容许伪度量：它满足距离的非负性、对称性和三角不等式，但不要求零距离的两个标号相同。

\begin{proposition}\label{g3:prop:finite-metric-matrices}
	固定整数 \(m\ge1\) 和 \(D\ge0\)。非空、至多含 \(m\) 个点且直径不超过 \(D\) 的度量空间，其等距类在 GH 距离下构成紧集。
\end{proposition}
\begin{proof}
	考虑所有满足
	\[
		a_{ii}=0,\qquad a_{ij}=a_{ji},\qquad
		0\le a_{ij}\le D,\qquad a_{ij}\le a_{i\ell}+a_{\ell j}
	\]
	的 \(m\times m\) 矩阵 \(A=(a_{ij})\)。这些条件定义了有限维闭方体 \([0,D]^{m^2}\) 的闭子集，因而矩阵集合紧。对每个 \(A\)，令 \(i\sim j\) 当且仅当 \(a_{ij}=0\)。三角不等式保证这确为等价关系，并且
	\[
		d_A([i],[j])=a_{ij}
	\]
	不依赖代表，成为商集合 \(X_A\) 上的距离。

	两个矩阵 \(A,B\) 给出对应
	\(\{([i]_A,[i]_B):1\le i\le m\}\)，其畸变不超过
	\(\max_{i,j}|a_{ij}-b_{ij}|\)。因此
	\begin{equation}\label{g3:eq:matrix-gh-bound}
		d_{\mathrm{GH}}(X_A,X_B)
		\le\frac12\max_{i,j}|a_{ij}-b_{ij}|.
	\end{equation}
	这说明矩阵到商空间等距类的映射连续。任意至多含 \(m\) 个点的空间都可以重复标号到 \(m\) 个，因而是某个 \(X_A\)。所述集合是紧矩阵集的连续像，故紧。
\end{proof}

这里没有把每个矩阵元单独取整。直接四舍五入可能破坏三角不等式；保留满足全部距离条件的矩阵，并在极限后合并零距离标号，才得到合法的度量空间。

\subsection{对角子列构造极限}
\label{g3:subsec:1003:02}

一层有限网只记录一个精度。为同时记住所有尺度，我们把各层的网按固定顺序编号，使粗网的标号在后续各层继续保留。由此得到可数张距离表，再对全部矩阵元作同一个对角抽取。极限标号给出的空间可能尚不完备，下面先明确如何补上其 Cauchy 极限。

\begin{lemma}\label{g3:lem:metric-completion}
	每个非空度量空间 \((S,d)\) 都有完备化：存在完备度量空间 \(\widehat S\) 及等距嵌入 \(S\hookrightarrow\widehat S\)，其像稠密。这样的完备化在保持 \(S\) 中各点的等距意义下唯一。
\end{lemma}
\begin{proof}
	取 \(S\) 中所有 Cauchy 序列，规定 \((s_j)\sim(t_j)\) 当且仅当 \(d(s_j,t_j)\to0\)，并令
	\[
		\widehat d([(s_j)],[(t_j)])=\lim_{j\to\infty}d(s_j,t_j).
	\]
	三角不等式给出
	\[
		|d(s_j,t_j)-d(s_\ell,t_\ell)|
		\le d(s_j,s_\ell)+d(t_j,t_\ell),
	\]
	所以右侧极限存在。再次使用三角不等式，可知换成等价序列不改变极限；所得函数满足距离公理，零距离正好对应同一等价类。常值序列把 \(S\) 等距嵌入这个商集合。

	若 \(z=[(s_j)]\)，常值序列所代表的点 \(s_j\) 在商距离下趋于 \(z\)：给定 \(\varepsilon>0\)，当 \(j,\ell\) 足够大时有 \(d(s_j,s_\ell)<\varepsilon\)，对 \(\ell\) 取极限即可。故 \(S\) 的像稠密。

	若 \((z_j)\) 是商空间中的 Cauchy 序列，利用稠密性选择 \(x_j\in S\)，使 \(\widehat d(x_j,z_j)<1/j\)。于是
	\[
		d(x_j,x_\ell)
		\le j^{-1}+\widehat d(z_j,z_\ell)+\ell^{-1},
	\]
	故 \((x_j)\) 在 \(S\) 中为 Cauchy 序列。其等价类 \(z=[(x_j)]\) 满足 \(x_j\to z\)，再由 \(\widehat d(x_j,z_j)\to0\) 得到 \(z_j\to z\)。这证明了完备性。

	最后，若 \(S\) 稠密等距地嵌入另一完备空间 \(Y\)，把 \([(s_j)]\) 送到 \((s_j)\) 在 \(Y\) 中的极限。等价性和距离连续性分别保证良定义与保距；稠密性保证满射。任何保持 \(S\) 的等距都须保持这些极限，因此它唯一。
\end{proof}

\begin{proposition}\label{g3:prop:uniform-nets-subsequence}
	设 \((X_i,d_i)\) 是非空紧度量空间序列。若存在 \(D<\infty\) 和对每个 \(k\ge1\) 给定的正整数 \(n_k\)，使
	\[
		\operatorname{diam}X_i\le D,
		\qquad N_{X_i}(2^{-k})\le n_k
		\qquad(i,k\ge1),
	\]
	则存在子列 \(X_{i_j}\) 和非空紧度量空间 \(X\)，满足 \(d_{\mathrm{GH}}(X_{i_j},X)\to0\)。极限可以由各层有限网的距离矩阵直接构造。
\end{proposition}
\begin{proof}
	证明分为三件事：先使标号间的距离同时收敛，再把极限距离补成紧空间，最后用同一层网控制 GH 误差。

	令 \(\varepsilon_k=2^{-k}\)、\(m_k=n_1+\cdots+n_k\)。在每个 \(X_i\) 中选取 \(\varepsilon_k\)-网 \(E_{i,k}\)，点数不超过 \(n_k\)。依次列出 \(E_{i,1},E_{i,2},\ldots\) 的点，每一层不足 \(n_k\) 个时重复其中一点补足，得到一列有标号的点
	\[
		p_i^1,p_i^2,\ldots\in X_i.
	\]
	前 \(m_k\) 个点中包含 \(E_{i,k}\)，因而形成 \(\varepsilon_k\)-网；粗层的标号也一直保留。允许重复使每个空间都有同样的标号数，不要求它们的网具有相同的组合形状。

	对每一对正整数 \(a,b\)，数列 \(d_i(p_i^a,p_i^b)\) 都落在紧区间 \([0,D]\) 内。把所有标号对排成一个序列，先为第一对选取收敛子列，再为第二对从中选取子列，逐次继续；取第 \(j\) 次子列的第 \(j\) 项得到对角子列。对任意固定的一对标号，该对角子列最终属于相应的收敛子列。因此沿同一个子列，所有极限
	\[
		d_\infty(a,b)=\lim_{j\to\infty}
		 d_{i_j}(p_{i_j}^a,p_{i_j}^b)
	\]
	都存在。以下把这个子列重新记为 \(i\)。非负性、对称性和三角不等式均可取极限，故 \(d_\infty\) 是 \(\mathbb N\) 上的伪度量。把零距离标号合并，得到度量空间 \(S\)；再用上一引理构造其完备化 \(X\)。以 \(p^a\in X\) 表示标号 \(a\) 的像。

	现在证明 \(X\) 紧。固定 \(k\) 及任意标号 \(a\)。原来的网给出
	\[
		\min_{1\le b\le m_k}d_i(p_i^a,p_i^b)\le\varepsilon_k.
	\]
	最小值只涉及有限个数，可以与极限交换，所以
	\[
		\min_{1\le b\le m_k}d_X(p^a,p^b)\le\varepsilon_k.
	\]
	因此稠密子集 \(S\) 被以 \(p^1,\ldots,p^{m_k}\) 为中心的有限个闭 \(\varepsilon_k\)-球覆盖。有限个闭球的并仍闭，它也包含 \(S\) 的闭包 \(X\)。这证明 \(X\) 全有界；它已按构造完备，故由\cref{g3:lem:complete-totally-bounded}紧。

	最后令
	\[
		F_i^k=\{p_i^a:1\le a\le m_k\},\qquad
		F^k=\{p^a:1\le a\le m_k\},
	\]
	都赋予各自空间的限制距离。带相同标号的点组成一个对应，即使有重复或零距离合并，其两个投影仍满射。该对应的畸变为至多
	\[
		\eta_{i,k}=\max_{1\le a,b\le m_k}
		\bigl|d_i(p_i^a,p_i^b)-d_X(p^a,p^b)\bigr|,
	\]
	而固定 \(k\) 时 \(\eta_{i,k}\to0\)。由有限网估计、对应公式及 GH 三角不等式，得到
	\begin{equation}\label{g3:eq:diagonal-gh-error}
	\begin{split}
		d_{\mathrm{GH}}(X_i,X)
		&\le d_{\mathrm{GH}}(X_i,F_i^k)
		+d_{\mathrm{GH}}(F_i^k,F^k)+d_{\mathrm{GH}}(F^k,X)\\
		&\le2\varepsilon_k+\tfrac12\eta_{i,k}.
	\end{split}
	\end{equation}
	给定误差目标，先把 \(k\) 取得足够大，再把 \(i\) 取得足够大，右侧即任意小。由此得到所需收敛。
\end{proof}

距离矩阵的逐项收敛本身只控制固定的标号。式\eqref{g3:eq:diagonal-gh-error}中的两项网误差把这个可数信息提升为对整个空间的统一控制；补全则把标号所逼近的全部极限点纳入空间。两个步骤各有作用。

\subsection{统一覆盖数的预紧性判据}
\label{g3:subsec:1003:03}

以 \(\mathcal M\) 表示非空紧度量空间的等距类所成的集合，赋予上一节已经建立的距离 \(d_{\mathrm{GH}}\)。紧空间有可数稠密子集，其等距类可以由可数距离矩阵的完备化编码，因此这里确实可以在一个集合上讨论度量。若 \(\mathcal F\subset\mathcal M\) 的闭包在 \(\mathcal M\) 中紧，就称 \(\mathcal F\) 在 GH 距离下\term[yu jin]{预紧}[precompact]。

\termidx[Gromov yujinxing dingli]{Gromov 预紧性定理}
\foreignidx{Gromov precompactness theorem}
\begin{theorem}[Gromov 预紧性判据]\label{g3:thm:gh-precompactness}
	设 \(\mathcal F\) 为一族非空紧度量空间的等距类。下列条件等价：
	\begin{equivalences}
		\item \(\mathcal F\) 在 GH 距离下预紧。
		\item 存在统一常数 \(D<\infty\)，使 \(\operatorname{diam}X\le D\) 对所有 \(X\in\mathcal F\) 成立；并且对每个 \(\varepsilon>0\)，存在有限整数 \(N(\varepsilon)\)，使
		\[
			N_X(\varepsilon)\le N(\varepsilon)
			\qquad(X\in\mathcal F).
		\]
		\item \(\mathcal F\) 作为 GH 度量空间的子集全有界。
	\end{equivalences}
\end{theorem}
\begin{proof}
	空族的结论直接成立，以下设 \(\mathcal F\ne\varnothing\)。预紧集的闭包紧，所以有每个尺度的有限网，这给出第一项蕴含第三项。

	设第三项成立。取一个有限 GH 距离的1-网，记网中空间为 \(Y_1,\ldots,Y_s\)。由\cref{g3:prop:gh-diameter}的直径估计，任意 \(X\in\mathcal F\) 都满足
	\[
		\operatorname{diam}X\le\max_j\operatorname{diam}Y_j+2.
	\]
	为了控制覆盖数，给定 \(\varepsilon>0\)，另取有限 GH 距离的 \(\varepsilon/8\)-网 \(Z_1,\ldots,Z_t\)。每个 \(X\) 与某个 \(Z_j\) 的 GH 距离小于 \(\varepsilon/4\)，故式\eqref{g3:eq:covering-number-transfer}给出
	\[
		N_X(\varepsilon)\le N_{Z_j}(\varepsilon/2)
		\le\max_{1\le j\le t}N_{Z_j}(\varepsilon/2)<\infty.
	\]
	有限个紧空间的覆盖数各自有限，所以右侧是统一界。这证明第二项。

	若第二项成立，任意 \(\mathcal F\) 中的序列都满足\cref{g3:prop:uniform-nets-subsequence}，从而有一个 GH 收敛到非空紧空间的子列。还须把这个序列结论用于闭包。取 \(Y_i\in\overline{\mathcal F}\)，为每个 \(i\) 选 \(X_i\in\mathcal F\) 使 \(d_{\mathrm{GH}}(X_i,Y_i)<1/i\)。某个子列 \(X_{i_j}\) 收敛到 \(X\in\mathcal M\)，三角不等式便使 \(Y_{i_j}\) 也收敛到 \(X\)。因为 \(\overline{\mathcal F}\) 闭，\(X\) 属于该闭包。故闭包顺序紧。它因而完备，因为每个 Cauchy 序列有收敛子列就整列收敛；它也全有界，否则可在某个固定正尺度逐次选出两两分离的点列，这种序列不可能有收敛子列。由\cref{g3:lem:complete-totally-bounded}，闭包紧，故第一项成立。
\end{proof}

\begin{example}\label{g3:exm:equilateral-nonprecompact}
	令 \(E_m\) 为含 \(m\) 个点、任意两不同点的距离都等于一的空间。它们的直径一致不超过一，但当 \(0<\varepsilon<1\) 时有 \(N_{E_m}(\varepsilon)=m\)，所以不预紧。更精确地，若 \(m>n\)，每个 \(E_m\) 与 \(E_n\) 间的对应都把某两个不同的 \(E_m\) 点配到同一个 \(E_n\) 点，畸变至少为一；任意满射产生的对应畸变又至多为一，故
	\[
		d_{\mathrm{GH}}(E_m,E_n)=\frac12.
	\]
	直径控制了整体尺度，却不能控制同一尺度上彼此分离的点有多少。另一方面，区间 \([0,m]\) 的直径趋于无穷，说明直径条件也有独立内容。
\end{example}

第八章的体积比较恰好能产生覆盖数上界。下面的应用也说明，在 GH 预紧性中，总体积不必有统一正下界。

\begin{corollary}\label{g3:cor:ricci-diameter-gh-precompact}
	固定整数 \(n\ge2\)、实数 \(k\) 及 \(D>0\)。所有满足
	\[
		\operatorname{Ric}_g\ge(n-1)kg,\qquad
		\operatorname{diam}(M,g)\le D
	\]
	的连通、紧、光滑无边界 \(n\) 维黎曼流形，在 GH 距离下预紧。
\end{corollary}
\begin{proof}
	由\cref{g3:cor:compact-riemannian-complete}，这些紧黎曼流形完备。置 \(k_0=\min\{k,0\}\)，则仍有 \(\operatorname{Ric}_g\ge(n-1)k_0g\)，而 \(k_0\) 的模型首零半径为无穷。给定 \(0<\varepsilon<D\)，逐个选取两两距离大于 \(\varepsilon\) 的点，直到不能再选。紧性保证过程有限终止：若产生无限点列，它没有 Cauchy 子列，与紧性矛盾。所得点 \(p_1,\ldots,p_N\) 构成闭 \(\varepsilon\)-网，开球 \(B(p_j,\varepsilon/2)\) 两两不交。

	因为 \(B(p_j,2D)=M\)，\cref{g3:thm:bishop-gromov}给出
	\[
		\operatorname{Vol}_g B(p_j,\varepsilon/2)
		\ge\frac{V_{k_0,n}(\varepsilon/2)}{V_{k_0,n}(2D)}
		\operatorname{Vol}_g(M).
	\]
	对互不相交的球求和，并消去正数 \(\operatorname{Vol}_g(M)\)，得到
	\begin{equation}\label{g3:eq:ricci-covering-number}
		N_M(\varepsilon)\le N
		\le\frac{V_{k_0,n}(2D)}{V_{k_0,n}(\varepsilon/2)}.
	\end{equation}
	右侧只依赖 \(n,k,D,\varepsilon\)。当 \(\varepsilon\ge D\) 时，一个球已可覆盖。统一直径界与各尺度的覆盖数界均已获得，故预紧性判据适用。
\end{proof}

体积在式\eqref{g3:eq:ricci-covering-number}中被约去，因而这个结论容许一个方向不断缩短的平坦环面。极限的维数与局部拓扑需要额外的几何信息来判断；下一节将直接计算这种塌缩。

\subsection{GH 距离的完备性与测地性的保持}
\label{g3:subsec:1003:04}

前面的构造已经给出了紧极限。现在把它用于 GH 距离自身的 Cauchy 序列，并检验长度几何能否随极限保留。这两件事分别控制空间序列的终点与极限中的最短线。

\begin{theorem}\label{g3:thm:gh-complete}
	非空紧度量空间的等距类空间 \((\mathcal M,d_{\mathrm{GH}})\) 完备。
\end{theorem}
\begin{proof}
	设 \((X_i)\) 为 GH 距离下的 Cauchy 序列。给定 \(\delta>0\)，可取 \(i_0\) 使所有 \(i\ge i_0\) 都满足 \(d_{\mathrm{GH}}(X_i,X_{i_0})<\delta\)；再把前面的有限个空间加入，就得到这列空间组成的集合的有限 \(\delta\)-网。因此这个集合在 GH 距离下全有界。

	由\cref{g3:thm:gh-precompactness}，某个子列 \(X_{i_j}\) 收敛到非空紧空间 \(X\)。任给 \(\varepsilon>0\)，Cauchy 条件保证当 \(i,\ell\) 足够大时有 \(d_{\mathrm{GH}}(X_i,X_\ell)<\varepsilon/2\)。再取足够大的 \(j\)，使 \(i_j\) 落在这个范围且 \(d_{\mathrm{GH}}(X_{i_j},X)<\varepsilon/2\)，就有
	\[
		d_{\mathrm{GH}}(X_i,X)
		\le d_{\mathrm{GH}}(X_i,X_{i_j})+d_{\mathrm{GH}}(X_{i_j},X)
		<\varepsilon.
	\]
	故整列收敛，证明完备性。
\end{proof}

同样的有限网方法还说明 \(\mathcal M\) 可分：所有有限度量空间在其中稠密，而对每个固定点数与直径界，\cref{g3:prop:finite-metric-matrices}给出的紧集有可数稠密子集。对点数和整数直径界取可数并，便得到 \(\mathcal M\) 的一个可数稠密子集。

\begin{proposition}\label{g3:prop:gh-geodesic-limit}
	设非空紧测地度量空间 \(X_i\) 在 GH 距离下收敛到非空紧度量空间 \(X\)。则 \(X\) 仍是测地空间。特别地，非空紧长度空间的 GH 极限仍为紧测地空间。
\end{proposition}
\begin{proof}
	选择对应 \(\mathcal R_i\subset X_i\times X\)，使 \(\operatorname{dis}\mathcal R_i\le\eta_i\to0\)。任取 \(x,y\in X\)，选择 \(x_i,y_i\in X_i\)，满足 \((x_i,x),(y_i,y)\in\mathcal R_i\)。在 \(X_i\) 中连接 \(x_i,y_i\) 的最短线有中点 \(z_i\)，即
	\[
		d_i(x_i,z_i)=d_i(y_i,z_i)=\tfrac12d_i(x_i,y_i).
	\]
	再取 \(w_i\in X\) 使 \((z_i,w_i)\in\mathcal R_i\)。畸变估计给出
	\[
		\bigl|d_X(x,w_i)-\tfrac12d_X(x,y)\bigr|\le\tfrac32\eta_i,
		\qquad
		\bigl|d_X(y,w_i)-\tfrac12d_X(x,y)\bigr|\le\tfrac32\eta_i.
	\]
	由于 \(X\) 紧，某个子列 \(w_{i_j}\) 收敛到 \(z\)。距离连续性使
	\[
		d_X(x,z)=d_X(y,z)=\tfrac12d_X(x,y).
	\]
	故任意两点都有精确中点。紧空间完备，\cref{g3:prop:midpoint-geodesic}便给出测地线：反复插入中点，在二进有理参数上定义保距的线段，再用完备性延拓到整个参数区间。

	最后，紧长度空间完备且局部紧，由\cref{g3:thm:metric-hopf-rinow}本身就是测地空间，所以第二句归入已经证明的情形。
\end{proof}

有限网逼近任意紧度量空间时，网本身通常没有非恒定的连续道路；这不会妨碍 GH 收敛。上一命题保留测地性的理由是能够把中点传到极限，再由完备性恢复整条线段。拓扑结构则有另一种行为：平坦环面的塌缩会改变维数，而测地性仍然保留。
% 纲要条目排为编号小节标题，并进入目录与书签；说明段仍保留于纲要注释中。

% !TeX root = ../../Book3.tex
% 纲要标题：### 10.4 非紧极限与切锥
% 纲要位置：几何学计划纲要.md，第 2484 行。
% 正文状态：初稿；尚未正式审核。

% 纲要内容开始
% 
% * 带基点收敛和球的控制
% * 重缩放及切锥
% * 平坦环面的塌缩
% * 切锥不唯一的反例
% 
% 纲要内容结束

\section{非紧极限与切锥}
\label{g3:sec:1004}

紧空间的 GH 距离同时比较整个空间。把越来越大的圆周从一个固定点附近观察，所见极限却是一条直线；远处的闭合部分逐渐离开每个有界观察范围。处理这样的极限，需要指定基点，并在每个固定半径内比较。放大一个点附近的几何也属于同一问题：原来半径趋零的球，在重缩放后成为固定半径的球。

\subsection{带基点收敛和球的控制}
\label{g3:subsec:1004:01}
先在紧空间上把基点也列为需要匹配的数据。对非空紧度量空间 \((A,p)\)、\((B,q)\)，定义
\begin{equation}\label{g3:eq:pointed-compact-gh}
	d_{\mathrm{GH}}^\bullet((A,p),(B,q))
	=\frac12\inf\{\operatorname{dis}\mathcal R:
	\mathcal R\subset A\times B\text{ 为对应},\ (p,q)\in\mathcal R\}.
\end{equation}
它也等于在公共度量空间中对
\(\max\{d_H(A,B),d(p,q)\}\) 取下确界。事实上，公共空间中距离小于 \(\varepsilon\) 的配对给出畸变至多为 \(2\varepsilon\) 的对应，并可包含基点对；反向使用\cref{g3:thm:gh-correspondence}证明中的桥接度量，每个配对点间的距离不超过所选桥长，基点也不例外。对应的复合保持基点对，所以三角不等式照常成立。在零距离的证明中，把基点放在稠密点列的第一位，所得等距必将 \(p\) 映到 \(q\)。因此式\eqref{g3:eq:pointed-compact-gh}在紧带基点空间的等距类上是度量。
\symidx{\(d_{\mathrm{GH}}^\bullet\)：紧带基点度量空间的 Gromov--Hausdorff 距离}

\begin{definition}\label{g3:def:pointed-gh-convergence}
	设 \((X_j,d_j,p_j)\) 与 \((X,d,p)\) 为非空、闭球紧的测地度量空间。若对每个 \(R>0\)，都有
	\begin{equation}\label{g3:eq:pointed-gh-balls}
		d_{\mathrm{GH}}^\bullet
		\bigl((\overline B_{X_j}(p_j,R),p_j),
		(\overline B_X(p,R),p)\bigr)\longrightarrow0,
	\end{equation}
	就称 \((X_j,p_j)\) \term[dai jidian GH shoulian]{带基点 GH 收敛}[pointed Gromov--Hausdorff convergence]于 \((X,p)\)，记为
	\((X_j,d_j,p_j)\xrightarrow{\mathrm{pGH}}(X,d,p)\)。所有闭球均取原空间距离的限制，不改取球内道路的内在距离。
\end{definition}

这里“闭球紧”采用\cref{g3:def:proper-metric-space}的约定。测地条件使半径略有偏差的点可以沿径向线移回球内；它正是下面球控制的依据。

\begin{lemma}\label{g3:lem:pointed-gh-ball-restriction}
	设 \((X,p)\)、\((Y,q)\) 为非空、闭球紧的测地度量空间。对 \(0<R\le S\)，有
	\[
		d_{\mathrm{GH}}^\bullet
		\bigl((\overline B_X(p,R),p),(\overline B_Y(q,R),q)\bigr)
		\le3d_{\mathrm{GH}}^\bullet
		\bigl((\overline B_X(p,S),p),(\overline B_Y(q,S),q)\bigr).
	\]
	因而在\cref{g3:def:pointed-gh-convergence}中，只检验正整数半径已经足够；也可检验一列趋于无穷的半径 \(R_j\)，要求相应两球的带基点 GH 距离趋于零。
\end{lemma}
\begin{proof}
	取半径 \(S\) 两球间的对应 \(\mathcal R\)，使 \((p,q)\in\mathcal R\) 且 \(\operatorname{dis}\mathcal R<\eta\)。若 \((x,y)\in\mathcal R\)，则
	\[
		|d_X(p,x)-d_Y(q,y)|<\eta.
	\]
	保留其中至少一个点属于半径 \(R\) 球的配对。将另一个点若超出半径 \(R\)，便沿从基点到它的一条最短线截到半径 \(R\)。每个配对至多移动一个点，移动量小于 \(\eta\)。所得关系仍覆盖两个小球，并包含基点对。比较两个新配对时，原距离误差小于 \(\eta\)，两次移动又至多增加 \(2\eta\)，故新畸变小于 \(3\eta\)。对 \(\eta\) 取下确界得到估计。

	整数半径的结论直接由此推出。若在半径 \(R_j\to\infty\) 上的误差趋零，固定任意 \(R\)，最终有 \(R<R_j\)，再用同一估计即可。反过来，若每个整数半径都收敛，选递增的整数 \(J_m\)，使 \(j\ge J_m\) 时半径 \(m\) 上的误差小于 \(1/m\)；在 \(J_m\le j<J_{m+1}\) 时取 \(R_j=m\)，便得到扩张球上的描述。
\end{proof}

对一般闭球紧空间，固定半径球可能在边界处跳变。例如 \(\{0,1+1/j\}\subset\mathbb R\) 收敛到 \(\{0,1\}\)，但以零为基点的半径一闭球始终只有一个点。一般版本的带基点收敛须允许半径误差或只使用适当的连续半径；本节的测地假设使式\eqref{g3:eq:pointed-gh-balls}成为可直接使用的描述。

\begin{proposition}\label{g3:prop:pointed-gh-unique}
	闭球紧测地空间的带基点 GH 极限在保持基点的等距意义下唯一。
\end{proposition}
\begin{proof}
	若同一序列收敛到 \((X,p)\)、\((Y,q)\)，紧带基点距离的三角不等式使每个整数 \(m\) 上的两球距离为零。因此存在保持基点的满等距
	\(F_m:\overline B_X(p,m)\to\overline B_Y(q,m)\)。这些映射未必彼此相容，下面用对角选择取得相容极限。

	闭球紧使 \(X\) 可分：对每个整数半径和每个正整数 \(l\)，选该球的有限 \(1/l\)-网，取这些网的并即可。取含 \(p\) 的可数稠密集 \(\{x_i\}\)。对固定 \(i\)，充分大的 \(m\) 都使 \(F_m(x_i)\) 有定义，并落在 \(Y\) 中半径 \(d_X(p,x_i)\) 的紧球内。逐点选子列，得到同一子列上全部 \(F_m(x_i)\) 收敛。极限保持所有距离，故由完备性唯一延拓为等距嵌入 \(F:X\to Y\)，且 \(F(p)=q\)。稠密点逼近及等距性还说明，沿所选子列，每个固定 \(x\in X\) 都满足 \(F_m(x)\to F(x)\)。

	给定 \(y\in Y\)，令 \(x_m=F_m^{-1}(y)\)，其中 \(m\) 充分大。它们位于 \(X\) 的同一个紧球中，取子列使 \(x_m\to x\)。于是
	\[
		d_Y(F(x),y)
		\le d_Y(F(x),F_m(x))+d_X(x,x_m)\longrightarrow0.
	\]
	因此 \(F(x)=y\)，得到满射性。
\end{proof}

非紧空间没有统一的整体直径界，但每个固定观察球仍可以有统一的有限网。上一节的距离矩阵方法可以同时记录全部整数半径上的点，从而构造一个共同的极限空间。

\begin{theorem}[带基点预紧性]\label{g3:thm:pointed-gh-precompactness}
	设 \((X_j,d_j,p_j)\) 是一列非空、闭球紧的测地度量空间。假设对每个 \(R>0\)、\(\varepsilon>0\)，存在与 \(j\) 无关的整数 \(N(R,\varepsilon)\)，使 \(\overline B_{X_j}(p_j,R)\) 有至多 \(N(R,\varepsilon)\) 个点的 \(\varepsilon\)-网。则存在带基点 GH 收敛子列，其极限也是闭球紧的测地度量空间。
\end{theorem}
\begin{proof}
	构造的关键是给所有半径中的网点统一编号，再一次性取所有点对距离的极限；不能把每个球独立得到的极限直接当成已有嵌套关系。

	对正整数 \(m,l\)，在每个 \(\overline B_{X_j}(p_j,m)\) 中选一个 \(2^{-l}\)-网，以重复基点的方式补足到固定的 \(N(m,2^{-l})\) 项。另把基点本身作为固定标签零。这样得到一个共同的可数标签集 \(I\)，标签 \(a\) 在 \(X_j\) 中所指的点记为 \(x_a^j\)，且每个标签都有与 \(j\) 无关的基点距离上界。每个点对的距离序列因而有界。按\cref{g3:prop:uniform-nets-subsequence}的对角方法选同一子列，使
	\[
		d_\infty(a,b)=\lim_j d_j(x_a^j,x_b^j)
	\]
	对全部 \(a,b\in I\) 存在。这个函数是伪度量；先把零距离标签合并，再按\cref{g3:lem:metric-completion}取完备化，得到 \((X,d)\)。基点零的像为 \(p\)，标签像构成稠密集 \(D\)。

	先证明闭球紧。固定 \(R\)，选整数 \(m>R+1\)。设 \(E_{m,l}\subset D\) 是半径 \(m\)、精度 \(2^{-l}\) 那组标签的极限像。若 \(x_a\in D\) 且 \(d(p,x_a)<m\)，则充分大的 \(j\) 有 \(x_a^j\in\overline B_{X_j}(p_j,m)\)。原网性质给出
	\[
		\min_{b\in I_{m,l}}d_j(x_a^j,x_b^j)\le2^{-l}.
	\]
	这里 \(I_{m,l}\) 是固定有限标签组，所以可以让有限个距离同时取极限，得到 \(d(x_a,E_{m,l})\le2^{-l}\)。再由 \(D\) 稠密，将此估计延拓到 \(\overline B_X(p,R)\)。这些覆盖球的中心未必属于所覆盖的闭球；从每个与 \(\overline B_X(p,R)\) 相交的覆盖球中选一个交点作为新中心，半径至多增至原来的两倍，就得到中心属于该闭球的有限网。因此每个闭球全有界；它又是完备空间中的闭集，由\cref{g3:lem:complete-totally-bounded}为紧集。

	再证明极限有测地线。先取 \(x_a,x_b\in D\)。在 \(X_j\) 中取 \(x_a^j,x_b^j\) 的最短线中点 \(z_j\)。这些中点距基点一致有界，故落在同一个整数半径 \(m\) 的球内。对每个 \(l\)，从该球的有限网选一点距 \(z_j\) 不超过 \(2^{-l}\)。有限性使某个标签 \(c_l\) 在无穷多个 \(j\) 上被选中，因此
	\[
		\left|d(x_a,x_{c_l})-\tfrac12d(x_a,x_b)\right|\le2^{-l},
		\qquad
		\left|d(x_b,x_{c_l})-\tfrac12d(x_a,x_b)\right|\le2^{-l}.
	\]
	全部距离的极限已固定，这里为确定一个标签而再取子列不会改变它们。点 \(x_{c_l}\) 落在同一个紧球内，取收敛子列就得到 \(x_a,x_b\) 的精确中点。一般 \(x,y\in X\) 可由 \(D\) 中点列逼近；相应中点仍在同一紧球中，取极限得到 \(x,y\) 的中点。由\cref{g3:prop:midpoint-geodesic}，完备性与精确中点保证 \(X\) 为测地空间。

	最后核对每个半径的收敛。固定 \(R>0\)，令 \(\delta=2^{-l}<1/4\)，仍取整数 \(m>R+1\)。在有限标签组 \(I_{m,l}\) 中保留满足 \(d(p,x_b)\le R+2\delta\) 的标签，并加入基点。对充分大的 \(j\)，它们的原像 \(x_b^j\) 距 \(p_j\) 不超过 \(R+3\delta\)。将这些点分别沿径向最短线截到半径 \(R\) 球，得到中心 \(y_b^j\)；在极限空间中作同样截断得到 \(y_b\)。相应移动量分别不超过 \(3\delta\)、\(2\delta\)。

	原空间的每个半径 \(R\) 球点都在某个原网点的 \(\delta\) 邻域内；该网点的基点距离不超过 \(R+\delta\)，故充分大的 \(j\) 时其标签属于所保留的组。因此 \(\{y_b^j\}\) 是 \(4\delta\)-网。前面已证 \(E_{m,l}\) 覆盖极限球，故 \(\{y_b\}\) 是 \(3\delta\)-网。有限距离矩阵收敛并计入截断移动量，给出
	\[
		\bigl|d_j(y_b^j,y_c^j)-d(y_b,y_c)\bigr|
		\le10\delta+o_j(1).
	\]
	把两球中靠近同标号中心的点配对，得到包含基点对的对应；其畸变至多
	\(2(4\delta)+2(3\delta)+10\delta+o_j(1)=24\delta+o_j(1)\)。因此两球的带基点 GH 距离上极限至多 \(12\delta\)。令 \(l\to\infty\)，便得到式\eqref{g3:eq:pointed-gh-balls}。
\end{proof}

基点固定以后，控制的是每个有界范围里的几何复杂程度。空间可以在更远处出现新的长枝、闭合或其他形状；只要这些特征不断远离基点，它们就不会出现在固定半径的极限球中。

\subsection{重缩放及切锥}
\label{g3:subsec:1004:02}
把一个长度为 \(r\) 的小物体放大到长度一，需要将距离乘以 \(r^{-1}\)。若原空间来自黎曼度量 \(g\)，相应的张量是 \(r^{-2}g\)，因为长度是度量平方值的平方根。

\begin{definition}\label{g3:def:metric-tangent-limit}
	设 \((X,d,p)\) 为非空、闭球紧的测地度量空间。若 \(r_j>0\)、\(r_j\to0\)，且
	\begin{equation}\label{g3:eq:metric-tangent-limit}
		(X,r_j^{-1}d,p)\xrightarrow{\mathrm{pGH}}(Y,d_Y,o),
	\end{equation}
	则称 \((Y,o)\) 为 \(X\) 在 \(p\) 处的一个\term[duliang qiejixian]{度量切极限}[metric tangent]，也称 GH 意义的\term[qiezhui]{切锥}[tangent cone]。这里的名称指重缩放极限；它是否具有真正的锥度量结构，以及是否与序列 \(r_j\) 无关，需要分别证明。
\end{definition}

由\cref{g3:thm:pointed-gh-precompactness}，若重缩放族的每个固定半径球具有统一的各尺度覆盖数，就能选出切极限。这个条件等价于在原空间中，对球 \(\overline B(p,Rr_j)\) 用半径 \(\varepsilon r_j\) 的有限网，并要求网点数只依赖 \(R,\varepsilon\)。因此切极限的存在控制的是相对尺度上的覆盖复杂度。

\begin{theorem}\label{g3:thm:smooth-riemannian-tangent}
	设 \((M^n,g)\) 为连通、完备、光滑无边界黎曼流形，\(n\ge1\)，\(p\in M\)。则
	\[
		(M,r^{-1}d_g,p)\xrightarrow[r\downarrow0]{\mathrm{pGH}}
		(\mathrm T_pM,|\cdot|_{g_p},0).
	\]
	特别地，光滑点的度量切极限唯一，且为 \(n\) 维 Euclid 空间。
\end{theorem}
\begin{proof}
	完备性和 Hopf--Rinow 定理使 \(M\) 的闭球紧。固定 \(R>0\)，在 \(p\) 处用正交单位基取法坐标。\Cref{g3:lem:uniform-normal-neighborhood}保证，当 \(r\) 足够小时，半径 \(3Rr\) 的闭距离球恰为相同半径闭切球的指数像，且该坐标在更大的邻域中光滑。定义
	\[
		F_r:\overline B_{g_p}(0,R)\longrightarrow
		\overline B_{r^{-1}d_g}(p,R),\qquad F_r(v)=\exp_p(rv).
	\]
	径向距离公式说明它是保持基点的双射。

	由\cref{g3:thm:normal-metric-expansion}，存在 \(\eta_r\to0\)，使在 \(|z|\le3Rr\) 上
	\[
		(1-\eta_r)|\xi|^2\le g_z(\xi,\xi)
		\le(1+\eta_r)|\xi|^2.
	\]
	设 \(|v|,|w|\le R\)。坐标直线段从 \(rv\) 到 \(rw\)，始终位于半径 \(Rr\) 的坐标球内，给出
	\[
		d_g(F_r(v),F_r(w))\le r\sqrt{1+\eta_r}\,|v-w|.
	\]
	下界必须兼顾离开坐标球的竞争道路。若道路始终留在半径 \(3Rr\) 的球中，逐段积分度量下界得到长度至少为 \(r\sqrt{1-\eta_r}|v-w|\)。若它离开该球，从初点到首次到达球面所付出的长度至少为
	\(3Rr-d_g(p,F_r(v))\ge2Rr\)，也不小于同一下界，因为 \(|v-w|\le2R\)。对全部道路取下确界，得到
	\[
		\sqrt{1-\eta_r}\,|v-w|
		\le r^{-1}d_g(F_r(v),F_r(w))
		\le\sqrt{1+\eta_r}\,|v-w|.
	\]
	故 \(F_r\) 的畸变一致趋于零。它的图是包含基点对的对应，式\eqref{g3:eq:pointed-compact-gh}说明每个半径 \(R\) 的两球收敛。由于这对所有 \(r\downarrow0\) 成立，所有趋零尺度序列都给出相同极限。
\end{proof}

另一类直接可算的切锥来自真正的度量锥。例如\cref{g3:exm:flat-metric-cone}中的圆锥 \(C_\theta\)，在顶点 \(o\) 的径向伸缩满足
\[
	d(\lambda x,\lambda y)=\lambda d(x,y)\qquad(\lambda>0).
\]
所以 \((C_\theta,r^{-1}d,o)\) 对每个 \(r>0\) 都与原锥保持基点等距，其顶点切锥就是自身。光滑点的极限是 Euclid 空间，锥顶的极限则可以保留角缺陷。

\subsection{平坦环面的塌缩}
\label{g3:subsec:1004:03}
令 \(S^1_L=\mathbb R/L\mathbb Z\) 为周长 \(L>0\) 的圆周，取弧长距离。考虑矩形平坦环面
\[
	T_{L,h}=S^1_L\times S^1_h,\qquad L,h>0,
\]
及乘积度量。第九章的平移群商把它实现为
\(\mathbb R^2/(L\mathbb Z\times h\mathbb Z)\)。两坐标分别选最近的格点代表，便得到
\begin{equation}\label{g3:eq:thin-torus-distance}
	d_{T_{L,h}}((x,y),(x',y'))^2
	=d_{S^1_L}(x,x')^2+d_{S^1_h}(y,y')^2.
\end{equation}

\begin{proposition}\label{g3:prop:flat-torus-collapse}
	固定 \(L>0\)。对于任意 \(h>0\)，投影 \(\pi:T_{L,h}\to S^1_L\) 为满射，且
	\[
		\operatorname{dis}\pi\le\frac h2,
		\qquad d_{\mathrm{GH}}(T_{L,h},S^1_L)\le\frac h4.
	\]
	因此当 \(h\downarrow0\) 时，二维平坦环面 GH 收敛到周长 \(L\) 的圆周。
\end{proposition}
\begin{proof}
	若两个圆周分量的距离分别为 \(a,b\)，则 \(0\le b\le h/2\)，并且
	\[
		0\le\sqrt{a^2+b^2}-a\le b\le h/2.
	\]
	这正是投影的畸变估计。投影图为对应，由\cref{g3:thm:gh-correspondence}得到 GH 上界。
\end{proof}

整个过程中截面曲率恒为零，而环面的面积为 \(Lh\)，趋于零。一个圆周方向的直径被压到零，极限中只剩另一个方向。这是可直接计算的\term[tasuo]{塌缩}[collapse]：距离收敛允许维数下降，环面的两条独立绕行也不会都保留为极限圆周的环路。

\ProseParagraphBreak

若一边塌缩、一边同时放大，结果取决于两个尺度的比值。这个现象可以完全由格的周期长度判断。

\begin{proposition}\label{g3:prop:torus-collapse-rescale}
	设 \(L>0\) 固定，\(h_j\downarrow0\)、\(r_j\downarrow0\)，基点均取 \(([0],[0])\)。则环面 \(T_{L,h_j}\) 的距离乘以 \(r_j^{-1}\) 后，有下列带基点极限：
	\[
	\begin{array}{c|c}
	 h_j/r_j\text{ 的极限}&\text{带基点极限}\\ \hline
	 0&\mathbb R\\
	 b\in(0,\infty)&\mathbb R\times S^1_b\\
	 +\infty&\mathbb R^2
	\end{array}
	\]
	右侧分别取标准距离、标准平坦乘积距离及原点基点。
\end{proposition}
\begin{proof}
	重缩放后的两个周期长度为
	\(A_j=L/r_j\to\infty\)、\(B_j=h_j/r_j\)。固定 \(R>0\)。当 \(A_j>4R\) 时，半径 \(R\) 球的水平坐标可唯一展开到 \([-R,R]\)，其中任意两点的水平差小于半周长；故该球连同限制距离与平坦圆柱 \(\mathbb R\times S^1_{B_j}\) 中的同半径球等距。

	若 \(B_j\to0\)，圆柱球投影到 \([-R,R]\) 为满射，畸变不超过 \(B_j/2\)，所以极限是直线的同半径球。若 \(B_j\to\infty\)，最终两个周期都大于 \(4R\)，两坐标均可展开，所得球与 Euclid 半径 \(R\) 闭圆盘等距。

	若 \(B_j\to b>0\)，用同一水平坐标和同一归一化圆周坐标 \(t\in\mathbb R/\mathbb Z\) 配对两圆柱。圆周距离分别为 \(B_j\min_{m\in\mathbb Z}|t-t'+m|\) 与 \(b\min_{m\in\mathbb Z}|t-t'+m|\)，差不超过 \(|B_j-b|/2\)。由乘积距离公式，全部配对的距离误差也不超过该数。选取至少一端在半径 \(R\) 球内的配对，将另一端沿径向最短线截回球内；与\cref{g3:lem:pointed-gh-ball-restriction}的证明相同，所得对应的畸变不超过 \(3|B_j-b|/2\)，且保留基点对。这证明了圆柱极限。
\end{proof}

对每个固定 \(h>0\)，先在环面的一点放大，得到 \(\mathbb R^2\)；先令 \(h\to0\) 得到圆周，再在圆周的一点放大，则得到 \(\mathbb R\)。因此改变取极限的次序会改变所观察的几何。上一命题的圆柱情形进一步记录了两个尺度相当时仍可见的闭合方向。

\subsection{切锥不唯一的反例}
\label{g3:subsec:1004:04}
光滑点的唯一切极限来自法坐标中度量系数的一致趋近。一般长度空间在越来越小的尺度中可以出现不同的几何构形。下面给出一个闭球紧测地空间，直接核算其两个不同切极限。

\begin{example}\label{g3:exm:nonunique-tangent-tree}
	令
	\[
		a_j=2^{-2^j}\qquad(j=1,2,\ldots).
	\]
	以半直线 \([0,\infty)\) 为主干，在主干点 \(a_j\) 处粘上一条长度 \(a_j\) 的闭线段。线段的一端与 \(a_j\) 识别，另一端自由；赋予所得空间 \(T\) 以沿树枝计算的道路距离，基点为主干原点 \(o=0\)。令 \(T_*\) 表示在半直线的点一处粘一条长度一闭线段的树，仍以主干原点为基点。那么
	\begin{equation}\label{g3:eq:nonunique-tree-tangents}
		(T,a_j^{-1}d,o)\xrightarrow{\mathrm{pGH}}(T_*,o),
		\qquad
		\bigl(T,(a_ja_{j+1})^{-1/2}d,o\bigr)
		\xrightarrow{\mathrm{pGH}}([0,\infty),0).
	\end{equation}
\end{example}
\begin{proof}
	任意两点之间都有唯一的简单道路：同一支段上的两点直接沿支段连接；不同支段上的点先回到各自附着点，再沿主干连接。这些道路的长度定义正定距离，并由相应道路实现，所以 \(T\) 是测地空间。它的闭球紧：对球中任意点列，若有子列位于主干及有限条支段的并中，则在一个有限紧图中有收敛子列；否则可选支段指标趋于无穷的子列，该子列上的点距 \(o\) 不超过 \(2a_j\to0\)，故收敛于 \(o\)。这也说明无限条支段在原点附近积聚，并没有造成缺失的极限点。

	先按 \(a_j\) 放大。第 \(j\) 条支段的附着距离及长度恰好都是一。比它更大的支段中，最近一条的附着距离为
	\[
		\frac{a_{j-1}}{a_j}\longrightarrow\infty;
	\]
	所以它们最终离开每个固定半径球。所有更小支段的长度至多
	\(\varepsilon_j=a_{j+1}/a_j\to0\)。把这些小支段收缩到其附着点，保留主干及第 \(j\) 条支段，得到到 \(T_*\) 的映射。在固定半径球上，当大支段已经离开时，该映射为满射并保持基点。每个点至多沿自己的支段移动 \(\varepsilon_j\)，所以任意一对点的距离误差不超过 \(2\varepsilon_j\)。式\eqref{g3:eq:pointed-compact-gh}给出第一个极限。

	再令 \(s_j=\sqrt{a_ja_{j+1}}\)。此时
	\[
		\frac{a_j}{s_j}=\sqrt{\frac{a_j}{a_{j+1}}}\longrightarrow\infty,
		\qquad
		\frac{a_{j+1}}{s_j}=\sqrt{\frac{a_{j+1}}{a_j}}\longrightarrow0.
	\]
	因此每个固定半径球内，较大的支段全部消失于远处，留下的支段长度一致趋零。将全部剩余支段投影到主干，在该球上为满射，畸变至多为剩余最大支长的两倍，得到第二个极限。
\end{proof}

\begin{figure}[htbp]
	\centering
	% !TeX root = ../Book3.tex
\begin{tikzpicture}[x=.95cm,y=.85cm,>=stealth]
	\draw[thick,->] (0,0)--(3,0);
	\fill (0,0) circle (1.6pt) node[below] {\(o\)};
	\foreach \x/\h in {.15/.15,.3/.3,.65/.65,1.5/1.3} {
		\draw[thick] (\x,0)--(\x,\h);
		\fill (\x,\h) circle (1.3pt);
	}
	\node[below] at (1.5,0) {\(a_j\)};
	\node[below] at (.65,-.3) {\(a_{j+1}\)};
	\node at (1.4,1.8) {原树 \(T\)};
	\draw[->] (3.1,.25)--(5.4,1.4) node[midway,above,sloped] {尺度 \(a_j\)};
	\draw[->] (3.1,-.25)--(5.4,-1.5) node[midway,below,sloped] {尺度 \(\sqrt{a_ja_{j+1}}\)};
	\begin{scope}[shift={(5.9,1.4)}]
		\draw[thick,->] (0,0)--(3,0);
		\draw[thick] (1,0)--(1,1);
		\fill (0,0) circle (1.6pt) node[below] {\(o\)};
		\fill (1,1) circle (1.3pt);
		\node[below] at (1,0) {\(1\)};
		\node[right] at (1,.6) {长 \(1\)};
		\node at (2,1.3) {\(T_*\)};
	\end{scope}
	\begin{scope}[shift={(5.9,-1.5)}]
		\draw[thick,->] (0,0)--(3,0);
		\fill (0,0) circle (1.6pt) node[below] {\(0\)};
		\node[above] at (1.5,0) {\([0,\infty)\)};
	\end{scope}
\end{tikzpicture}

	\caption{同一棵树的两种切极限。左图的支段位置与长度仅作示意；按 \(a_j\) 放大时保留距离一处的单位支段，按相邻尺度的几何均值放大时只剩半直线。}
	\label{g3:fig:tangent-scales}
\end{figure}
如\cref{g3:fig:tangent-scales}所示，放大时被保留的支段取决于所选尺度。两个极限不等距：\(T_*\) 在距基点一处有分枝点，删除该点后出现三个连通分支；半直线上的任何点删除后至多出现两个。\(T_*\) 也不是以 \(o\) 为顶点的度量锥，因为真正的度量锥有任意正比例的径向相似变换；这会把距离一处的分枝点变成任意距离处的分枝点，而 \(T_*\) 只有一个分枝点。

\ProseParagraphBreak

这个例子区分了三个问题：重缩放族是否存在收敛子列，各子列是否给出同一个带基点空间，以及所得空间是否具有锥结构。它发生在明确给出的测地树上。关于某一类带曲率界流形的奇异极限，还须另外利用该类的曲率、维数与非塌缩条件；光滑点的唯一 Euclid 切极限则已经由\cref{g3:thm:smooth-riemannian-tangent}证明。
% 纲要条目排为编号小节标题，并进入目录与书签；说明段仍保留于纲要注释中。

% !TeX root = ../../Book3.tex
% 纲要标题：### 10.5* 合成曲率的基本模型
% 纲要位置：几何学计划纲要.md，第 2491 行。
% 正文状态：初稿；尚未正式审核。

% 纲要内容开始
% 
% * CAT 空间的比较定义
% * Alexandrov 下界及三角条件
% * 光滑曲率与合成条件的连接
% 
% ---
% 
% 纲要内容结束

\OptionalSection{合成曲率的基本模型}{g3:sec:1005}

第九章把截面曲率转成了三角形距离的不等式。这些不等式只比较有限个点，因而在没有切空间的树、凸集或度量极限上仍能提出。本节用可直接计算的模型说明上、下比较的差异，再回到光滑情形，证明小三角形比较确实能够恢复截面曲率的界。

\subsection{CAT 空间的比较定义}
\label{g3:subsec:1005:01}
沿用\cref{g3:def:cat-comparison}：测地空间 \(X\) 为 \(\mathrm{CAT}(k)\)，是指每个周长小于 \(2b_k\) 的最短三角形，与常曲率 \(k\) 模型中同边长的三角形相比，任意两个对应边点的距离都不增大。这里
\[
	b_k=\begin{cases}\pi/\sqrt{k},&k>0,\\\infty,&k\le0.\end{cases}
\]
每一种最短边的选择都要检验。局部 \(\mathrm{CAT}(k)\) 的意思是每点有一个邻域，其内在度量使之成为 \(\mathrm{CAT}(k)\) 空间；邻域内在距离与全空间距离的关系须在具体使用时核对。

\begin{proposition}\label{g3:prop:cat-zero-midpoint-criterion}
	设 \((X,d)\) 为测地空间。它为 \(\mathrm{CAT}(0)\)，当且仅当对每条测地线段 \([x,y]\) 的中点 \(m\) 和每个 \(z\in X\)，都有
	\begin{equation}\label{g3:eq:cat-midpoint-inequality}
		d(z,m)^2\le\frac{d(z,x)^2+d(z,y)^2}{2}-\frac{d(x,y)^2}{4}.
	\end{equation}
	此时测地线段唯一，而且对其常速参数化 \(\gamma:[0,1]\to X\)，有
	\begin{equation}\label{g3:eq:metric-cn-inequality}
		d(z,\gamma(t))^2\le(1-t)d(z,x)^2+td(z,y)^2-t(1-t)d(x,y)^2.
	\end{equation}
\end{proposition}
\begin{proof}
	在 Euclid 比较三角形中，中线平方满足式\eqref{g3:eq:cat-midpoint-inequality}右端的等式，CAT 上比较便给出必要性。

	反过来，沿给定线段令 \(D=d(x,y)\)、\(F(t)=d(z,\gamma(t))^2\)。将中点条件用于任意子段，说明连续函数 \(F(t)-D^2t^2\) 满足中点凸性。反复二等分得到二进分数处的凸性不等式，再由连续性延拓到所有参数，便得式\eqref{g3:eq:metric-cn-inequality}。这首先给出顶点到对边的上比较。

	还需得到两个边点间的比较。对三角形 \(pqr\)，取 \(u\in[p,q]\)、\(v=\gamma_{p,r}(t)\)。先对顶点 \(r\) 与边点 \(u\) 使用上式，得 \(d(r,u)^2\le|\bar r-\bar u|^2\)。再以 \(u\) 为固定点，对边 \([p,r]\) 使用上式，得到
	\[
	\begin{aligned}
		d(u,v)^2
		&\le(1-t)d(u,p)^2+t d(u,r)^2-t(1-t)d(p,r)^2\\
		&\le(1-t)|\bar u-\bar p|^2+t|\bar u-\bar r|^2
		-t(1-t)|\bar p-\bar r|^2
		=|\bar u-\bar v|^2.
	\end{aligned}
	\]
	不同两边总有共同顶点，同一边上的对应距离相等，故全部 CAT 条件成立。

	若 \(m,m'\) 是同端点两条线段的中点，在式\eqref{g3:eq:cat-midpoint-inequality}中取 \(z=m'\)，右侧为零，所以 \(m=m'\)。再对两半段依次应用同一论证，两条线段在所有二进参数处相同，连续性使它们处处相同。
\end{proof}

\begin{example}\label{g3:exm:real-tree-cat-zero}
	测地空间中若任意两点之间只有一条简单弧，且这条弧就是连接它们的测地线段，就称它为\term[shi shu]{实树}[real tree]，也记作 \(\mathbb R\)-树。有限树图赋予正边长及道路距离便是例子。三个点的连接线段组成三脚架，允许某条支腿长度为零。

	每棵实树都是 \(\mathrm{CAT}(0)\)。为验证中点判据，把 \([x,y]\) 按弧长记为 \(\gamma:[0,L]\to X\)。从 \(z\) 到该段的分支汇合点为 \(b=\gamma(a)\)，设 \(h=d(z,b)\)。由树的唯一弧性质，
	\[
		d(z,\gamma(t))=h+|t-a|.
	\]
	于是
	\[
		d(z,\gamma(t))^2-t^2=h^2+a^2-2at+2h|t-a|
	\]
	为凸函数。这正给出式\eqref{g3:eq:metric-cn-inequality}，从而给出中点判据。树上的分叉与 CAT 上界完全相容。
\end{example}

闭凸 Euclid 子集也是 \(\mathrm{CAT}(0)\)：最短段为直线段，任意三个点及其边位于一个至多二维的仿射平面中，模型比较处处取等号。它可能有边界或角点。另一方面，带长度度量的圆周不是 \(\mathrm{CAT}(0)\)，因为对径点有两条最短半圆；这与刚证明的唯一性矛盾，但每段足够小的弧都是 \(\mathrm{CAT}(0)\)。

\begin{proposition}\label{g3:prop:compact-cat-zero-limit}
	若非空紧 \(\mathrm{CAT}(0)\) 测地空间 \(X_j\) 在 GH 意义下收敛到非空紧度量空间 \(X\)，则 \(X\) 也是 \(\mathrm{CAT}(0)\)。
\end{proposition}
\begin{proof}
	\Cref{g3:prop:gh-geodesic-limit}保证 \(X\) 测地。取 \(x,y,z\in X\) 和任意选定的中点 \(m\)。用畸变趋零的对应取 \(x_j,y_j,z_j,m'_j\)，使四点间距离趋于相应极限距离。在 \(X_j\) 中另取 \(x_j,y_j\) 的真正中点 \(m_j\)。将中点不等式用于顶点 \(m'_j\)，得到
	\[
		d_j(m'_j,m_j)^2\le\frac{d_j(m'_j,x_j)^2+d_j(m'_j,y_j)^2}{2}
		-\frac{d_j(x_j,y_j)^2}{4}\longrightarrow0.
	\]
	因此 \(d_j(z_j,m_j)\to d(z,m)\)。再将 \(X_j\) 的中点不等式用于 \(z_j\) 并取极限，就得到 \(X\) 的中点不等式。上一判据适用，结论成立。
\end{proof}

\subsection{Alexandrov 下界及三角条件}
\label{g3:subsec:1005:02}
下界使用相反的距离方向，但不要求最短线唯一。继续采用同边长模型，对边上的对应点按相同弧长比例选取。

\begin{definition}\label{g3:def:cbb-triangle-condition}
	取 \(k\in\mathbb R\)。记 \(s_k\) 为 \(s_k''+ks_k=0\)、\(s_k(0)=0\)、\(s_k'(0)=1\) 的解，\(b_k=\pi/\sqrt k\) 当 \(k>0\)，否则 \(b_k=\infty\)。
	设 \((X,d)\) 为测地空间，\(k\in\mathbb R\)。若对任意选定的最短三角形 \(pqr\)，只要周长小于 \(2b_k\)，每个顶点 \(p\) 与对边上的点 \(z\) 就满足
	\begin{equation}\label{g3:eq:cbb-triangle-condition}
		d(p,z)\ge d_k(\bar p,\bar z),
	\end{equation}
	其中 \(d_k\) 为常曲率 \(k\) 的二维单连通完备模型距离，\(\bar p,\bar z\) 为同边长比较三角形中的对应点。则称它满足 \(\mathrm{CBB}(k)\) 三角条件，即具有\term[Alexandrov xia qulv jie]{Alexandrov 下曲率界}[Alexandrov lower curvature bound] \(k\)。这里 CBB 为 curvature bounded below 的缩写。本节采用顶点到对边的表述，并明确允许退化三角形；完备性、局部紧性等要求在各个结果中另行列出。
\end{definition}

当 \(k=0\) 时，式\eqref{g3:eq:cbb-triangle-condition}等价于沿每条常速最短段的平方距离下比较
\begin{equation}\label{g3:eq:cbb-zero-distance}
	d(p,\gamma_{q,r}(t))^2
	\ge(1-t)d(p,q)^2+td(p,r)^2-t(1-t)d(q,r)^2.
\end{equation}
这是 Euclid 模型中平方距离的精确公式。在 \(k\ne0\) 时，可用第九章的 \(m_k\) 将模型距离换成满足 \(h''+kh=1\) 的函数。它将三角条件与一维比较方程联系起来。

\begin{example}\label{g3:exm:tripod-no-lower-bound}
	取 \(k\in\mathbb R\)。记 \(s_k\) 为 \(s_k''+ks_k=0\)、\(s_k(0)=0\)、\(s_k'(0)=1\) 的解，\(c_k=s_k'\)、\(\operatorname{ct}_k=c_k/s_k\)，\(b_k=\pi/\sqrt k\) 当 \(k>0\)，否则 \(b_k=\infty\)。
	三条长 \(\varepsilon\) 的线段共用一个端点 \(o\)，另三个端点记为 \(p,q,r\)。它们的三边长都是 \(2\varepsilon\)，而 \(o\) 是边 \([q,r]\) 的中点。实际距离为 \(d(p,o)=\varepsilon\)，Euclid 模型中线长为 \(\sqrt3\varepsilon\)，所以三脚架不满足 \(\mathrm{CBB}(0)\)。

	更强地，在一个真正分叉点附近，实树不满足任何固定有限 \(k\) 的局部下曲率界。取足够小的 \(\varepsilon\)，三条支腿仍可等长截取，且周长 \(6\varepsilon<2b_k\)。模型函数 \(s_k(t)=t+O(t^3)\)、\(c_k(t)=1-kt^2/2+O(t^4)\) 代入余弦公式，说明模型中线为 \(\sqrt3\varepsilon+O(\varepsilon^3)\)，仍严格大于 \(\varepsilon\)。当 \(k=0\) 时为上面的精确值。因此局部 CAT 上比较并不排除分叉，下比较则会检测它。
\end{example}

闭凸 Euclid 子集同时满足上、下零曲率比较。正曲率球面在周长小于首个共轭尺度两倍的范围内与球面模型完全相同；它的对径点有多条最短线，也不妨碍下比较。第九章还已证明每个连通完备、光滑无边界且 \(K\ge k\) 的黎曼流形满足式\eqref{g3:eq:cbb-triangle-condition}。这包括平坦环面，虽然它不是全局 \(\mathrm{CAT}(0)\)。

\begin{table}[htbp]
	\centering
	\caption{三个零曲率比较模型；下比较取本节的顶点到对边形式。}
	\label{g3:tab:synthetic-models}
	\begin{tabularx}{\textwidth}{@{}lXX@{}}
		\toprule
		空间 & CAT(0) 上比较 & CBB(0) 下比较\\
		\midrule
		闭凸 Euclid 子集 & 成立，模型等号 & 成立，模型等号\\
		有分叉的实树 & 成立 & 分叉点任意小尺度失败\\
		二维及以上平坦环面 & 局部成立，整体失败 & 整体成立\\
		\bottomrule
	\end{tabularx}
\end{table}
表中所有判断都由已经给出的距离公式、最短线唯一性或第九章的比较定理得到。它们说明，合成曲率条件能够区别具有同样局部平坦部分的不同粘合方式。

\subsection{光滑曲率与合成条件的连接}
\label{g3:subsec:1005:03}
光滑流形上的比较可以双向解释：Hessian 比较给出有限三角形的不等式，而将对边缩到一个点，再将径向距离缩到零，能够从这些不等式恢复截面曲率。下面把这个过程写出。

\begin{theorem}\label{g3:thm:smooth-synthetic-curvature}
	取 \(k\in\mathbb R\)。记 \(s_k\) 为 \(s_k''+ks_k=0\)、\(s_k(0)=0\)、\(s_k'(0)=1\) 的解，\(b_k=\pi/\sqrt k\) 当 \(k>0\)，否则 \(b_k=\infty\)。
	设 \((M^n,g)\) 连通、完备、光滑且无边界，\(n\ge2\)、\(k\in\mathbb R\)。则：
	\begin{enumerate}
		\item \(K\le k\) 当且仅当每点都有强凸邻域，使其内在度量为 \(\mathrm{CAT}(k)\)。
		\item \(K\ge k\) 当且仅当每点都有强凸邻域，使其中每个足够小的最短三角形满足 \(\mathrm{CBB}(k)\) 条件。此时整个 \(M\) 满足周长小于 \(2b_k\) 的下比较。
	\end{enumerate}
\end{theorem}
\begin{proof}
	先证明曲率上界给局部上比较。由\cref{g3:thm:strongly-convex-neighborhood}，可选足够小的强凸邻域 \(U\)，其两点间距离平方光滑，且在 \(k>0\) 时使 \(\operatorname{diam}U<b_k/3\)。强凸性保证邻域内在距离就是环境距离，所用径向线均处于相应起点的割点之前。

	若 \(K\le k\)，\cref{g3:thm:hessian-comparison}给出
	\[
		\operatorname{Hess}(m_k\circ d(p,\cdot))
		\ge c_k(d(p,\cdot))g.
	\]
	在中心这个复合函数仍光滑，其 Hessian 等于 \(g\)。沿任意单位速对边，令 \(h=m_k(d(p,\sigma(t)))\)，便有 \(h''+kh\ge1\)。同边长模型中的 \(\bar h\) 满足等式且端点值相同。\Cref{g3:lem:one-dimensional-support-gluing}证明中的正权最小值论证应用于 \(\bar h-h\)，给出 \(h\le\bar h\)。这里边长小于 \(b_k\)，所有实际和模型距离也小于 \(b_k\)，所以 \(m_k\) 递增，得到顶点到对边的上比较。

	为了从顶点比较得到所有边点比较，记 \(H_{A,B}\) 为区间 \([0,L]\) 上满足 \(H''+kH=1\)、端点为 \(A,B\) 的解。两个这样的解之差满足齐次方程，因而
	\[
		H_{A',B'}(t)-H_{A,B}(t)
		=\frac{s_k(L-t)}{s_k(L)}(A'-A)
		+\frac{s_k(t)}{s_k(L)}(B'-B).
	\]
	当 \(L<b_k\) 时，两个系数非负，所以边值解随端点值增加。现在取三角形 \(pqr\) 中的 \(u\in[p,q]\)、\(v\in[p,r]\)。先比较 \(r\) 到 \(u\)，得
	\(m_k(d(r,u))\le m_k(d_k(\bar r,\bar u))\)；再以 \(u\) 为基点沿 \([p,r]\) 比较，并使用刚证明的端点单调性。端点 \(p\) 处的距离相等，端点 \(r\) 处为上述上界，故
	\[
		m_k(d(u,v))\le m_k(d_k(\bar u,\bar v)).
	\]
	再由单调性得到完整的 CAT 上比较。两次模型边值解确有同一几何解释：原比较三角形中的 \(\bar u,\bar p,\bar r\) 仍在同一模型内，子三角形周长不超过原周长。

	若 \(K\ge k\)，则\cref{g3:thm:toponogov-triangle}直接给出整个 \(M\) 的下比较，也给出任意强凸小邻域的相同条件。这证明两个正向结论。

	现在从局部三角比较恢复曲率。固定 \(p\)，沿单位速径向线取很近的 \(q=\exp_p(rv)\)，再取过 \(q\) 的单位速测地线 \(\sigma(t)\)。令 \(h(t)=m_k(d(p,\sigma(t)))\)。对充分小的 \(\varepsilon>0\)，三角形的对边取 \(\sigma|_{[-\varepsilon,\varepsilon]}\)。模型距离函数 \(\bar h\) 满足 \(\bar h''+k\bar h=1\)，由偶、奇基本解可得
	\[
		\bar h(0)=\frac{h(-\varepsilon)+h(\varepsilon)-2m_k(\varepsilon)}{2c_k(\varepsilon)}.
	\]
	局部上比较说 \(h(0)\le\bar h(0)\)，下比较则反向。用
	\(c_k(\varepsilon)=1-k\varepsilon^2/2+O(\varepsilon^4)\)、
	\(m_k(\varepsilon)=\varepsilon^2/2+O(\varepsilon^4)\)
	展开并除以 \(\varepsilon^2\)，得到上比较时 \(h''(0)+kh(0)\ge1\)，下比较时 \(\le1\)。

	令 \(\dot\sigma(0)\perp\dot\gamma(r)\)。此时一阶径向导数为零，故上述不等式分别成为
	\[
		\operatorname{Hess}d(p,\cdot)(\dot\sigma(0),\dot\sigma(0))
		\ge\operatorname{ct}_k(r),\qquad\text{或}\qquad\le\operatorname{ct}_k(r).
	\]
	在平行识别下，\cref{g3:prop:radial-riccati,g3:prop:distance-hessian-jacobi}给出
	\[
		\operatorname{Hess}d(p,\cdot)(w_r,w_r)
		=\frac1r-\frac r3\langle R(w,v)v,w\rangle+O(r^2),
		\qquad\operatorname{ct}_k(r)=\frac1r-\frac{kr}{3}+O(r^3),
	\]
	其中 \(v,w\in\mathrm T_pM\) 正交且单位，\(w_r\) 为平行移动。减去 \(1/r\)，除以 \(r\) 后令 \(r\downarrow0\)，上比较推出 \(K(v,w)\le k\)，下比较推出 \(K(v,w)\ge k\)。任取点与二维平面即得两个逆向结论。
\end{proof}

当 \(k=0\) 时，\cref{g3:thm:hadamard-cat-zero}进一步说明：完备、单连通的非正截面曲率流形是全局 \(\mathrm{CAT}(0)\)。若去掉单连通性，平坦环面表明只能从局部上比较出发。下界的全局化则已由 Toponogov 完成，不需要单连通条件。

\ProseParagraphBreak

这些模型也解释了为什么曲率比较适合研究极限。紧 CAT(0) 空间的极限仍满足中点不等式，却可以失去光滑坐标；实树甚至可以在一个点附近有不同尺度的支枝。对后续几何分析而言，度量极限先提供一个可比较的底空间。要在极限上保留微分方程、体积或光滑结构，还需要对应的额外控制与收敛概念。
% 纲要条目排为编号小节标题，并进入目录与书签；说明段仍保留于纲要注释中。


\begin{chaptersummary}
曲线长度把一个度量空间中的可行道路与端点距离联系起来。完备、局部紧的长度空间具有紧闭球，近乎最短曲线的中点可以取极限，再以二进参数补成测地线。锥、等距群商和闭子集粘合使这一套语言同时覆盖光滑与非光滑模型。

\ProseParagraphBreak

GH 距离等于对应最小畸变的一半，因而可以用有限距离表控制。紧空间的零 GH 距离给出满射等距，统一直径与各尺度覆盖数则给出预紧性。极限的构造依次使用对角子列、零距离商与完备化；这些步骤也证明 GH 等距类空间完备，并说明紧测地空间的极限仍有测地线。

\ProseParagraphBreak

带基点收敛把整体比较改为每个固定球上的比较。光滑点的度量切极限是唯一的 Euclid 空间，塌缩环面在不同尺度比下却可趋向直线、圆柱或平面。稀疏支树进一步展示了切极限可以不唯一，也可以不具有锥结构。有限三角比较在这些空间上仍有意义：实树满足 CAT(0)，分叉却破坏任何固定的局部下曲率界。下一编将回到函数与微分方程，在具体分析假设下研究几何估计与紧性。
\end{chaptersummary}

\BookExercises
\begin{exercise}\label{g3:ex:snowflake-length}
	取 \(0<\alpha<1\)，在 \(\mathbb R\) 上令 \(d_\alpha(x,y)=|x-y|^\alpha\)。证明它是完备、闭球紧的度量，但每条非恒定连续曲线都有无穷长度。因此其内在距离在不同点间为 \(+\infty\)。
\end{exercise}
\begin{solution}
	对 \(u,v\ge0\)，有 \((u+v)^\alpha\le u^\alpha+v^\alpha\)：当 \(u+v>0\) 时，令 \(t=u/(u+v)\)，利用 \(t^\alpha\ge t\)、\((1-t)^\alpha\ge1-t\) 即得。它与通常的三角不等式一起证明 \(d_\alpha\) 的三角不等式，其余度量公理直接成立。\(d_\alpha\) 与原距离有相同的收敛列和 Cauchy 列，故完备；其闭半径 \(R\) 球是普通半径 \(R^{1/\alpha}\) 闭区间，所以紧。

	若连续曲线 \(\gamma\) 非恒定，限制到某子区间并必要时改变方向，可设其两端值为 \(a<b\)。对正整数 \(N\)，用介值定理依次取首次到达 \(a+j(b-a)/N\) 的时刻，得到一个递增的分割。因此
	\[
		L_{d_\alpha}(\gamma)\ge N\left(\frac{b-a}{N}\right)^\alpha
		=(b-a)^\alpha N^{1-\alpha}\longrightarrow\infty.
	\]
	连接不同点的连续曲线必非恒定，故内在距离为无穷。闭球紧本身不能替代长度空间假设。
\end{solution}

\begin{exercise}\label{g3:ex:circle-reflection-quotient}
	设 \(S^1_L=\mathbb R/L\mathbb Z\) 带弧长距离，\(L>0\)，反射 \(\iota([t])=[-t]\) 生成二阶等距群。证明其度量商与区间 \([0,L/2]\) 等距，给出显式等距映射，并解释商中两个端点来自何处。
\end{exercise}
\begin{solution}
	每个轨道恰有一个代表 \(s\in[0,L/2]\)，故映射 \([[t]])\mapsto d_{S^1_L}([0],[t])\) 是到该区间的双射。对两个代表 \(s,t\in[0,L/2]\)，商距离为
	\[
		\min\{|s-t|,s+t,L-s-t\}=|s-t|.
	\]
	这里 \(s+t\ge|s-t|\)；若 \(s\ge t\)，则 \(L-s-t\ge s-t\) 等价于 \(L\ge2s\)，另一情形同理。因此该双射等距。区间端点对应圆周上两个反射不动点 \([0]\)、\([L/2]\)。即使原圆周没有边界，非自由等距作用的商也可以产生边界点。
\end{solution}

\begin{exercise}\label{g3:ex:max-norm-geodesics}
	在 \(\mathbb R^2\) 上取 \(d_\infty((x,y),(x',y'))=\max\{|x-x'|,|y-y'|\}\)。证明这个空间完备、局部紧且为测地空间，并刻画从 \((0,0)\) 到 \((1,0)\)、参数区间为 \([0,1]\) 的全部常速测地线段。
\end{exercise}
\begin{solution}
	此范数与 Euclid 范数满足 \(\|v\|_\infty\le|v|\le\sqrt2\|v\|_\infty\)，因此空间完备、闭球紧。任意仿射线段按常速参数化都实现范数距离，所以它为测地空间。

	若 \(\gamma(t)=(x(t),y(t))\) 是所问线段，则端点距离为1，因而 \(d_\infty(\gamma(t),(0,0))=t\)、\(d_\infty(\gamma(t),(1,0))=1-t\)。前式给 \(x(t)\le t\)，后式给 \(x(t)\ge t\)，所以 \(x(t)=t\)。其第二分量必须为1-Lipschitz 函数，且 \(y(0)=y(1)=0\)。反过来，任何满足这些条件的 \(y\) 都给出测地线，因为
	\[
		d_\infty(\gamma(s),\gamma(t))=
		\max\{|s-t|,|y(s)-y(t)|\}=|s-t|.
	\]
	例如 \(y(t)=c\min\{t,1-t\}\)、\(|c|\le1\)，给出无穷多条不同线段。长度空间的完备性与局部紧性保证存在性，不保证唯一性。
\end{solution}

\begin{exercise}\label{g3:ex:gh-equilateral-spaces}
	设 \(E_n(a)\) 是由 \(n\ge2\) 个点组成的度量空间，其中任意两个不同点的距离均为 \(a>0\)。
	\begin{enumerate}
		\item 若 \(n=m\ge2\)，证明
		\[
			d_{\mathrm{GH}}(E_n(a),E_n(b))=\tfrac12|a-b|.
		\]
		\item 若 \(n>m\ge2\)，证明
		\[
			d_{\mathrm{GH}}(E_n(a),E_m(b))
			=\tfrac12\max\{a,|a-b|\}.
		\]
		\item 给定任意正数列 \((a_n)_{n\ge2}\)，证明 \(E_n(a_n)\) 在 \(n\to\infty\) 时 GH 收敛到某个非空紧度量空间，当且仅当 \(a_n\to0\)；并确定极限。
	\end{enumerate}
\end{exercise}
\begin{solution}
	1. 任取两边点集间的双射，其图像对应的畸变为 \(|a-b|\)，所以对应公式给出上界。两空间的直径分别为 \(a,b\)，直径下界给出同一个数，故取等号。

	2. 先证明下界。任取对应 \(R\)，对 \(E_n(a)\) 的每个点选一个配对的靶点。由于 \(n>m\)，至少两个不同源点 \(x,x'\) 被配到同一靶点 \(y\)，于是
	\[
		\operatorname{dis}R\ge|d(x,x')-d(y,y)|=a.
	\]
	直径差还给出 \(\operatorname{dis}R\ge|a-b|\)。因此所求 GH 距离至少为题给右侧。

	为得到上界，任取满射 \(f:E_n(a)\to E_m(b)\)，以其图像作对应。两个源点相同则误差为零；两个源点不同而像相同，则误差为 \(a\)；两个像不同，则误差为 \(|a-b|\)。由于 \(n>m\ge2\)，后两种情形都出现，故图像畸变恰为 \(\max\{a,|a-b|\}\)。对应公式给出上界。

	3. 若 \(a_n\to0\)，单点公式给出 \(d_{\mathrm{GH}}(E_n(a_n),\{*\})=a_n/2\to0\)，所以极限为单点空间。反过来，若这列空间 GH 收敛，则它是 GH-Cauchy 列。给定 \(\varepsilon>0\)，可取 \(N\) 使 \(n>m\ge N\) 时距离小于 \(\varepsilon\)。第2问的公式给出 \(a_n/2\le d_{\mathrm{GH}}(E_n(a_n),E_m(a_m))<\varepsilon\)。固定 \(m=N\)，便知所有 \(n>N\) 都满足 \(a_n<2\varepsilon\)。因此 \(a_n\to0\)，其极限等距于单点空间。
\end{solution}

\begin{exercise}\label{g3:ex:gh-product-stability}
	设 \(X,Y,A,B\) 为非空紧度量空间。在乘积上采用平方和距离，例如
	\[
		d_{X\times A}((x,a),(x',a'))
		=\sqrt{d_X(x,x')^2+d_A(a,a')^2}.
	\]
	证明这些乘积仍为紧度量空间，并证明
	\[
		d_{\mathrm{GH}}(X\times A,Y\times B)
		\le\sqrt{d_{\mathrm{GH}}(X,Y)^2+d_{\mathrm{GH}}(A,B)^2}.
	\]
	推导 \(d_{\mathrm{GH}}(X\times A,X)\le\operatorname{diam}A/2\)，并说明这个常数在哪种情形达到。
\end{exercise}
\begin{solution}
	平方和距离的三角不等式由两个因子的三角不等式和 \(\mathbb R^2\) 中的 Minkowski 不等式得到，其余度量公理直接成立。对乘积中的任意序列，先在第一个紧因子取收敛子列，再在第二个紧因子取收敛子列，得到乘积中收敛的子列；度量空间中的序列紧性等价于紧性，所以乘积紧。

	取 \(R\subset X\times Y\)、\(S\subset A\times B\) 为对应。定义乘积对应
	\[
		T=\{((x,a),(y,b)):(x,y)\in R,\ (a,b)\in S\}.
	\]
	两个投影都满射。由 Euclid 范数的反三角不等式，对两对这样的点有
	\[
	\begin{aligned}
		&\left|\sqrt{d_X(x,x')^2+d_A(a,a')^2}
		-\sqrt{d_Y(y,y')^2+d_B(b,b')^2}\right|\\
		&\qquad\le\sqrt{|d_X(x,x')-d_Y(y,y')|^2
		+|d_A(a,a')-d_B(b,b')|^2}\\
		&\qquad\le\sqrt{(\operatorname{dis}R)^2+(\operatorname{dis}S)^2}.
	\end{aligned}
	\]
	因此 \(\operatorname{dis}T\) 满足同一上界。对给定 \(\delta>0\)，分别取 \(R,S\) 使其畸变小于 \(2d_{\mathrm{GH}}(X,Y)+\delta\)、\(2d_{\mathrm{GH}}(A,B)+\delta\)。应用对应公式并令 \(\delta\downarrow0\)，即得乘积不等式。

	令 \(Y=X\)、\(B=\{*\}\)，并将 \(X\times\{*\}\) 与 \(X\) 等距识别，再用单点距离公式，得到所求推论。当 \(X\) 本身是单点时，乘积与 \(A\) 等距，这个上界就是单点公式中的等号。因此常数 \(1/2\) 不能统一减小。
\end{solution}

\begin{exercise}\label{g3:ex:gh-dense-subsets}
	设 \(X,Y\) 为非空紧度量空间，\(A\subset X\)、\(B\subset Y\) 为非空稠密子集，赋予限制距离。对这些未必紧的子集，暂用公共等距嵌入及同一 Hausdorff 下确界公式定义 \(\widehat d_{\mathrm{GH}}(A,B)\)。证明
	\[
		\widehat d_{\mathrm{GH}}(A,B)=d_{\mathrm{GH}}(X,Y).
	\]
	由此说明：若从本节的零距离定理中删掉紧性，并仍使用同一公式，零距离未必推出两个原空间等距。
\end{exercise}
\begin{solution}
	先注意，在任意度量空间中，非空有界集合的点集距离满足 \(d(z,A)=d(z,\overline A)\)；它又是连续函数，所以对另一个集合取上确界时，把这个集合换成闭包也不改变结果。因此
	\[
		d_H(A,B)=d_H(\overline A,\overline B)
	\]
	于任意公共环境中成立，只要按相同公式解释有界集合的 Hausdorff 距离。

	将紧空间 \(X,Y\) 等距嵌入任意公共空间，\(A,B\) 的像的闭包就是 \(X,Y\) 的像。于是这次嵌入给出的两个 Hausdorff 距离相等，取下确界得到 \(\widehat d_{\mathrm{GH}}(A,B)\le d_{\mathrm{GH}}(X,Y)\)。

	反过来，给定 \(A,B\) 在度量空间 \(Z\) 中的等距副本，用\cref{g3:lem:metric-completion}将 \(Z\) 等距放入其完备化 \(\overline Z\)。嵌入 \(A\to\overline Z\) 唯一延拓为等距嵌入 \(X\to\overline Z\)：对 \(x\in X\) 取 \(a_j\in A\) 使 \(a_j\to x\)，其像是 Cauchy 序列，定义延拓值为其极限；等距性保证极限与所取序列无关并保留任意点对的距离。\(B\to\overline Z\) 同样延拓到 \(Y\)。这些延拓的像分别为原像的闭包，所以其 Hausdorff 距离与原来 \(A,B\) 的相同。


	对原来的公共嵌入取下确界，得到 \(d_{\mathrm{GH}}(X,Y)\le\widehat d_{\mathrm{GH}}(A,B)\)，从而等式成立。最后取 \(X=Y=[0,1]\)、\(A=\mathbb Q\cap[0,1]\)、\(B=[0,1]\)。两者的扩展距离为零，但 \(A\) 可数、\(B\) 不可数，故不可能等距。公共嵌入的距离在这里识别了它们的完备化。
\end{solution}

\begin{exercise}\label{g3:ex:shrinking-star-gh-limit}
	把区间 \([0,1/j]\)、\(j=1,2,\ldots\) 的左端点全部识别为一个点 \(o\)。所得集合记为 \(T\)，其中第 \(j\) 条支上坐标为 \(t>0\) 的点记为 \((j,t)\)，赋予距离
	\[
		d(o,(j,t))=t,\qquad
		d((i,s),(j,t))=
		\begin{cases}|s-t|,&i=j,\\s+t,&i\ne j.\end{cases}
	\]
	令 \(T_m\) 只保留前 \(m\) 条支。
	\begin{enumerate}
		\item 证明这些公式定义测地距离，且 \(T\) 紧。
		\item 证明 \(d_{\mathrm{GH}}(T_m,T)\le1/(m+1)\)。
		\item 证明 \(o\) 的任意邻域都含无穷多条完整的支，并据此判断 \(T\) 是否为拓扑流形。
	\end{enumerate}
\end{exercise}
\begin{solution}
	1. 非负性、对称性及零距离条件直接成立。两点同支时，第三点在同支便应用实直线的三角不等式；第三点在异支时，两段距离之和为 \(s+t+2u\ge|s-t|\)。两点异支时，第三点若在其中一支，则所需不等式归为 \(|s-u|+u\ge s\)；若在第三支，则两段距离之和为 \(s+t+2u\ge s+t\)。因此三角不等式成立。同支点沿区间连接，异支点先到 \(o\) 再沿另一支前进，这些路径的长度恰等于所给距离，故 \(T\) 为测地空间。

	给定 \(\varepsilon>0\)，所有长度不超过 \(\varepsilon\) 的支都落在 \(\overline B(o,\varepsilon)\) 内，其余支只有有限条，每条紧区间可取有限 \(\varepsilon\)-网，故 \(T\) 全有界。为证完备，设 \((x_i)\) 是 Cauchy 序列。如果其中有一个子列落在有限条支的并内，则该有限并紧，子列有极限，而 Cauchy 条件使整列趋于同一极限。如果没有这样的子列，则各支的编号趋于无穷；相应坐标不超过支长，故 \(d(o,x_i)\to0\)。两种情形都收敛，所以 \(T\) 完备，进而紧。

	2. 将 \(T_m,T\) 置于公共空间 \(T\)。前者是后者的子集，任何删去的支上的点到 \(o\in T_m\) 的距离不超过 \(1/(m+1)\)，所以
	\[
		d_{\mathrm{GH}}(T_m,T)\le d_H^T(T_m,T)\le\frac1{m+1}.
	\]

	3. 任意含 \(o\) 的开邻域都含某个 \(B(o,r)\)，长度小于 \(r\) 的支全部包含其中，且这样的支有无穷多条。删去 \(o\) 后，不同支之间不能由连续道路连接：每条支去掉 \(o\) 后在 \(T\setminus\{o\}\) 中既开又闭。若 \(o\) 有流形坐标邻域，取其中对应于小球或边界半球的更小邻域，则删去 \(o\) 后至多有两个连通分支：一维内部点为两个，一维边界点为一个，维数至少二时连通。实际邻域删去 \(o\) 后却有无穷多个非空互相分离的支，因此 \(o\) 不是任何有限维拓扑流形的点。零维情形也不成立，因为 \(o\) 不是孤立点。
\end{solution}

\begin{exercise}\label{g3:ex:doubling-gh-precompact}
	设 \(C\ge1\) 为整数、\(D>0\)。称度量空间具有 \(C\)-倍增覆盖性质，是指每个闭半径 \(r>0\) 球都可由至多 \(C\) 个中心在空间内、半径为 \(r/2\) 的闭球覆盖。设一族非空紧度量空间 \(X\) 都具有这个性质，且 \(\operatorname{diam}X\le D\)。
	\begin{enumerate}
		\item 证明当 \(0<\varepsilon<D\) 时，
		\[
			N_X(\varepsilon)\le C^{\lceil\log_2(D/\varepsilon)\rceil},
		\]
		并推出这一族 GH 预紧。
		\item 若这些空间中的一列 \(X_i\) GH 收敛到紧空间 \(X\)，证明 \(X\) 也具有 \(C\)-倍增覆盖性质。
	\end{enumerate}
\end{exercise}
\begin{solution}
	1. 任取 \(p\in X\)，则 \(X=\overline B(p,D)\)。对覆盖中的每个球反复使用倍增性质，\(j\) 次以后得到至多 \(C^j\) 个半径为 \(D/2^j\) 的闭球。取 \(j=\lceil\log_2(D/\varepsilon)\rceil\)，便有 \(D/2^j\le\varepsilon\)，所以这些球的中心形成所需的有限网。\(\varepsilon\ge D\) 时一个中心已足够。所有尺度都有统一覆盖数界，Gromov 预紧性判据因而适用。

	2. 取畸变不超过 \(\eta_i\to0\) 的对应 \(\mathcal R_i\subset X_i\times X\)。固定 \(x\in X\)、\(r>0\)，选择 \((x_i,x)\in\mathcal R_i\)。在 \(X_i\) 中，用至多 \(C\) 个半径为 \((r+\eta_i)/2\) 的闭球覆盖 \(\overline B(x_i,r+\eta_i)\)，其中心记为 \(a_i^1,\ldots,a_i^C\)，不足 \(C\) 个时重复一个中心。选 \(b_i^j\in X\) 使 \((a_i^j,b_i^j)\in\mathcal R_i\)。紧性及有限次子列抽取使同一个子列上每个 \(b_i^j\) 都收敛，记极限为 \(b^j\)。

	任取 \(y\in\overline B(x,r)\)，选择 \((y_i,y)\in\mathcal R_i\)，则 \(d_i(x_i,y_i)\le r+\eta_i\)。它落在上述某个球内，因而
	\[
		\min_{1\le j\le C}d_X(y,b_i^j)
		\le\frac{r+\eta_i}{2}+\eta_i.
	\]
	对这个已经选定的子列取极限，得到 \(\min_jd_X(y,b^j)\le r/2\)。有限个极限中心不依赖 \(y\)，故它们的闭 \(r/2\)-球覆盖整个 \(\overline B(x,r)\)，证明倍增覆盖性质得以保留。
\end{solution}

\begin{exercise}\label{g3:ex:gh-connected-limit}
	设非空紧连通度量空间 \(X_i\) 在 GH 距离下收敛到非空紧度量空间 \(X\)。证明 \(X\) 连通。再给出每个空间恰有两个连通分支、极限却为一个区间的序列，并估计其 GH 收敛误差。
\end{exercise}
\begin{solution}
	若 \(X\) 不连通，可写成两个不交非空闭开子集的并 \(X=A\sqcup B\)。由于 \(A,B\) 紧且不交，连续距离函数在 \(A\times B\) 上取得正最小值
	\[
		\delta=d_X(A,B)>0.
	\]
	取充分大的 \(i\)，使存在畸变 \(\eta_i<\delta\) 的对应 \(\mathcal R_i\subset X_i\times X\)。令 \(A_i\) 为所有与某个 \(A\) 中的点对应的 \(X_i\) 点，\(B_i\) 类似定义。两个投影满射保证它们非空且并为 \(X_i\)。它们不交：若一个点同时与 \(a\in A,b\in B\) 对应，畸变便至少为 \(d_X(a,b)\ge\delta\)，产生矛盾。

	对 \(u\in A_i,v\in B_i\)，选对应的 \(a\in A,b\in B\)，得到
	\[
		d_i(u,v)\ge d_X(a,b)-\eta_i\ge\delta-\eta_i>0.
	\]
	所以 \(A_i,B_i\) 是相距正距离的两个非空集合，各自在 \(X_i\) 中开，构成一个分离，与 \(X_i\) 连通矛盾。这个证明通过对应直接搬运正距离的分离，没有要求近似映射连续。

	最后，在实直线中取 \(Y_m=[-1,-1/m]\cup[1/m,1]\)，\(m\ge2\)，赋予限制距离。每个 \(Y_m\) 恰有两个连通分支，而
	\[
		d_{\mathrm{GH}}(Y_m,[-1,1])
		\le d_H^{\mathbb R}(Y_m,[-1,1])=\frac1m\longrightarrow0.
	\]
	两个分支之间的间隙随 \(m\) 缩小，在极限中消失。因此极限连通不能反推逼近空间连通。
\end{solution}

\begin{exercise}\label{g3:ex:moving-basepoint-intervals}
	在区间 \(X_j=[-j,j]\) 上取通常距离。分别选取基点 \(p_j=0\)、\(q_j=j\)，以及 \(z_j=j-a\)，其中 \(a\ge0\) 固定，且只考虑 \(j>a\)。求三个带基点 GH 极限，并说明为何同一列紧空间 \(X_j\) 本身没有紧 GH 收敛子列。
\end{exercise}
\begin{solution}
	基点为零时，对每个固定 \(R>0\)，一旦 \(j>R\)，半径 \(R\) 闭球就是 \([-R,R]\)。因此 \((X_j,0)\) 带基点收敛到 \((\mathbb R,0)\)。

	基点为 \(q_j=j\) 时，用坐标 \(u=j-x\) 将 \((X_j,q_j)\) 等距识别为 \(([0,2j],0)\)。对固定 \(R\)，最终的半径 \(R\) 球就是 \([0,R]\)，所以极限为 \(([0,\infty),0)\)。

	基点为 \(z_j=j-a\) 时，用 \(u=x-(j-a)\) 得到区间 \([-2j+a,a]\)，基点为零。固定半径球最终等于 \([-R,\min\{a,R\}]\)，故极限为 \(((-\infty,a],0)\)。当 \(a>0\) 时，基点到唯一端点的距离为 \(a\)；它与端点作为基点的半直线不保持基点等距。当 \(a=0\) 时，反向参数化将此极限等距识别为第二种。

	另一方面，\(\operatorname{diam}X_j=2j\to\infty\)。由\cref{g3:prop:gh-diameter}，紧 GH 收敛序列的直径必须收敛到极限的有限直径，所以这里不存在紧 GH 收敛子列。带基点收敛只比较固定半径的球，允许整体直径发散。
\end{solution}

\begin{exercise}\label{g3:ex:three-scale-flat-tori}
	记 \(S^1_a=\mathbb R/a\mathbb Z\) 为周长 \(a>0\) 的平坦圆周。令 \(0<\varepsilon<1\)，在三维平坦环面
	\[
		T_\varepsilon=S^1_1\times S^1_\varepsilon\times S^1_{\varepsilon^2}
	\]
	上取乘积度量，基点为三个圆周的零点。固定 \(\alpha>0\)，把距离乘以 \(\varepsilon^{-\alpha}\)。分别求
	\(0<\alpha<1\)、\(\alpha=1\)、\(1<\alpha<2\)、\(\alpha=2\)、\(\alpha>2\)
	时的带基点 GH 极限，并给出各情形在固定半径球上的误差理由。
\end{exercise}
\begin{solution}
	重缩放后的三个圆周周期长度为
	\[
		\varepsilon^{-\alpha},\qquad
		\varepsilon^{1-\alpha},\qquad
		\varepsilon^{2-\alpha}.
	\]
	趋于无穷的周期在固定半径球上展开为直线；趋于零的周期可以投影掉；恒等于一的周期则保留为 \(S^1_1\)。因此五种极限依次为
	\[
		\mathbb R,\qquad
		\mathbb R\times S^1_1,\qquad
		\mathbb R^2,\qquad
		\mathbb R^2\times S^1_1,\qquad
		\mathbb R^3.
	\]
	各空间均取标准乘积距离及零点基点。

	具体核对如下。固定 \(R>0\)，所有趋于无穷的周期最终大于 \(4R\)。原点半径 \(R\) 球中的这些坐标可展开到 \([-R,R]\)，任意两点的对应差不超过 \(2R\)，故不会取到较短的绕行距离；这些坐标上的距离已经与 Euclid 距离完全相同。

	设需要投影掉的圆周周期为 \(b_1,\ldots,b_s\)。乘积中它们的联合直径为
	\[
		D_\varepsilon=\frac12\sqrt{b_1^2+\cdots+b_s^2}\longrightarrow0.
	\]
	将这些坐标置零，得到从半径 \(R\) 球到目标半径 \(R\) 球的满射：每个目标点都可用零纤维坐标提升。若保留坐标的两点距离为 \(u\)，丢弃坐标的距离为 \(v\le D_\varepsilon\)，则总距离为 \(\sqrt{u^2+v^2}\)，所以投影畸变至多为 \(D_\varepsilon\)。这给出带基点 GH 误差至多为 \(D_\varepsilon/2\)。当 \(\alpha>2\) 时没有需要丢弃的坐标，固定半径球最终已经与 Euclid 三维闭球等距。
\end{solution}

\begin{exercise}\label{g3:ex:all-tangents-of-sparse-tree}
	令 \(a_j=2^{-2^j}\)，在主干 \([0,\infty)\) 的点 \(a_j\) 处各粘上一条长 \(a_j\) 的闭支段，赋予所得树 \(T\) 道路距离，以主干原点 \(o=0\) 为基点；这是\cref{g3:exm:nonunique-tangent-tree}的稀疏支树。对 \(c>0\)，以 \(T_c\) 表示在半直线距原点 \(c\) 处粘上一条长度 \(c\) 支段的树。证明 \(T\) 在 \(o\) 处的全部带基点切极限恰为
	\[
		([0,\infty),0),\qquad (T_c,0)\quad(c>0),
	\]
	并证明这些空间两两不保持基点等距。
\end{exercise}
\begin{solution}
	正文已经证明半直线出现。对固定 \(c>0\)，取尺度 \(r_j=a_j/c\)。第 \(j\) 条支段放大后恰在距离 \(c\) 处且长 \(c\)；更大的支段离开每个固定球，更小的支段最大长度为 \(ca_{j+1}/a_j\to0\)。将小支段投影到主干，畸变至多为此最大长度的两倍，于是极限是 \(T_c\)。

	反过来，设某个趋零尺度序列 \(r_i\) 已给出带基点极限 \(Y\)。选整数 \(j_i\)，使
	\[
		a_{j_i+1}<r_i\le a_{j_i}.
	\]
	则 \(j_i\to\infty\)。记
	\[
		q_i=\frac{a_{j_i}}{r_i}\ge1,\qquad
		t_i=\frac{a_{j_i+1}}{r_i}\in(0,1),\qquad
		\frac{t_i}{q_i}=\frac{a_{j_i+1}}{a_{j_i}}\longrightarrow0.
	\]
	再取子列，使 \(q_i\to q\in[1,\infty]\)、\(t_i\to t\in[0,1]\)。其他更大支段的最近附着距离满足
	\[
		\frac{a_{j_i-1}}{r_i}
		=\frac{a_{j_i-1}}{a_{j_i}}q_i\longrightarrow\infty,
	\]
	而其他更小支段的最大长度满足
	\[
		\frac{a_{j_i+2}}{r_i}
		=\frac{a_{j_i+2}}{a_{j_i+1}}t_i\longrightarrow0.
	\]
	因此在每个固定观察球内，至多只有指标 \(j_i\)、\(j_i+1\) 的两条支段可能留下非零尺度。

	若 \(q<\infty\)，则 \(t=0\)，唯一保留的支段位置与长度均趋于 \(q\)，极限为 \(T_q\)。若 \(q=\infty\) 且 \(t>0\)，则第一条支段离开有界球，第二条趋于位置和长度均为 \(t\) 的支段，极限为 \(T_t\)。若 \(q=\infty\) 且 \(t=0\)，两条支段分别远离或缩小，极限是半直线。

	这里支段参数 \(c_i\to c>0\) 时 \(T_{c_i}\to T_c\) 的事实也可直接验证：把所有主干及支段的弧长参数乘以 \(c/c_i\)，得到两树间的满双射，它将全部距离同乘 \(c/c_i\)。在半径 \(R\) 球上，其距离误差至多为 \(2R|c/c_i-1|\)；将半径略超出的点沿径向线截回后，仍得到趋零畸变的对应。结合前述小支段投影，就完成了三种子列极限的核对。原序列已经收敛到 \(Y\)，由\cref{g3:prop:pointed-gh-unique}，它与所列相应极限保持基点等距。

	最后，\(T_c\) 恰有一个删除后产生三个连通分支的点，其到基点的距离为 \(c\)。保持基点的等距必须保持这个点及该距离，所以 \(c\ne c'\) 时 \(T_c\) 与 \(T_{c'}\) 不保持基点等距。半直线没有这样的分枝点，故与任何 \(T_c\) 都不等距。
\end{solution}

\begin{optionalexercise}\label{g3:ex:cat-norm-parallelogram}
	设实赋范向量空间 \((V,\|\cdot\|)\) 的距离满足 \(\mathrm{CAT}(0)\) 比较。证明其范数满足平行四边形恒等式
	\[
		\|x+y\|^2+\|x-y\|^2=2\|x\|^2+2\|y\|^2.
	\]
	据此证明 \(\mathbb R^2\) 上的 \(\ell^p\) 范数在 \(1\le p\le\infty\)、\(p\ne2\) 时不满足 \(\mathrm{CAT}(0)\)。
\end{optionalexercise}
\begin{solution}
	赋范空间中的仿射线段是测地线。把中点不等式用于端点 \(x,y\)、中点 \((x+y)/2\) 及顶点零，得到
	\[
		\|x+y\|^2+\|x-y\|^2\le2\|x\|^2+2\|y\|^2.
	\]
	再把这条不等式用于 \(u=x+y\)、\(v=x-y\)，并用范数的齐次性，得到
	\[
		4\|x\|^2+4\|y\|^2\le2\|x+y\|^2+2\|x-y\|^2.
	\]
	两式合并即为恒等式。

	对 \(x=(1,0)\)、\(y=(0,1)\)，当 \(p<\infty\) 时，恒等式要求 \(2\cdot2^{2/p}=4\)，恰在 \(p=2\) 时成立。\(p=\infty\) 时左侧为2、右侧为4，也不成立。因此其余这些范数均不满足 CAT(0)。其中某些范数仍严格凸并具有唯一最短段，所以测地唯一性只是 CAT 条件的一项后果。
\end{solution}

\begin{optionalexercise}\label{g3:ex:cat-product-collapse}
	设 \(X,Y\) 是 \(\mathrm{CAT}(0)\) 测地空间。证明平方和乘积距离
	\[
		d((x,y),(x',y'))=\sqrt{d_X(x,x')^2+d_Y(y,y')^2}
	\]
	使 \(X\times Y\) 仍为 \(\mathrm{CAT}(0)\)。若 \(X,Y\) 非空紧，说明将第二因子距离乘以 \(\varepsilon>0\) 后的空间 GH 收敛到 \(X\)，并给出误差界。
\end{optionalexercise}
\begin{solution}
	分别在两因子取同一参数区间 \([0,1]\) 的常速最短段，其乘积满足测地距离等式，故乘积测地。每个乘积中点的两个分量都是相应因子的中点。确实，若端点各分量距离组成向量 \(D=(D_1,D_2)\)，中点到两端的分量距离组成非负向量 \(a,b\)，则
	\[
		D_i\le a_i+b_i,\qquad |a|=|b|=|D|/2.
	\]
	于是 \(|D|\le|a+b|\le|a|+|b|=|D|\)，全部取等号。第一处等号使各分量 \(D_i=a_i+b_i\)，第二处等号使 \(a,b\) 同向；模相等迫使 \(a=b=D/2\)，结论得证。

	将两个因子的中点不等式相加，恰得到乘积距离平方的中点不等式。由\cref{g3:prop:cat-zero-midpoint-criterion}，乘积为 CAT(0)。将度量乘常数不改变零曲率比较，故缩放后的乘积也为 CAT(0)。到 \(X\) 的投影满射，且其畸变不超过 \(\varepsilon\operatorname{diam}Y\)，因此
	\[
		d_{\mathrm{GH}}(X\times\varepsilon Y,X)
		\le\tfrac12\varepsilon\operatorname{diam}Y\longrightarrow0.
	\]
	这里 \(\varepsilon Y\) 表示距离已乘 \(\varepsilon\) 的空间，不是一个新底集合。
\end{solution}

\begin{optionalexercise}\label{g3:ex:circle-cbb-zero}
	用圆周的距离公式直接证明：周长 \(L>0\) 的长度圆周满足 \(\mathrm{CBB}(0)\) 三角条件。允许三点位于半圆内，也允许所选三条最短边绕完整个圆周。结合正文，说明这个空间的上、下零曲率比较有何不同。
\end{optionalexercise}
\begin{solution}
	取最短边 \(qr\)，长度为 \(a>0\)，按弧长 \(t\in[0,a]\) 参数化。若顶点 \(p\) 在该弧内部，则三点沿这条长度至多 \(L/2\) 的弧排列，所有相关距离就是区间距离，模型三角形退化，比较取等号。

	否则，圆周中从 \(p\) 到弧内点的路径必须先经过 \(q\) 或 \(r\)。令 \(c=d(p,q)\)、\(b=d(p,r)\)，于是
	\[
		\rho(t)=d(p,\gamma(t))=\min\{c+t,b+a-t\}.
	\]
	这个等式也可由圆周上的两条弧逐一计算；即使到某个端点的最短弧曾经过另一端点，所得候选路径也不会减小另一条真实连接的最小值。Euclid 比较三角形中的对应距离平方为
	\[
		Q(t)=\left(1-\frac ta\right)c^2+\frac ta b^2-t(a-t).
	\]
	直接相减得到
	\[
	\begin{aligned}
		(c+t)^2-Q(t)&=\frac ta\bigl((a+c)^2-b^2\bigr)\ge0,\\
		(b+a-t)^2-Q(t)&=\frac{a-t}{a}\bigl((a+b)^2-c^2\bigr)\ge0,
	\end{aligned}
	\]
	其中使用了三角不等式。取两式左侧候选值的最小，得到 \(\rho(t)^2\ge Q(t)\)，即所需下比较。顶点与端点重合的情形由退化模型直接成立。每种最短边选择都可照此处理，故圆周满足 CBB(0)。它局部 CAT(0)，但对径点的两条最短半圆排除了整体 CAT(0)；下比较容许这样的整体识别。
\end{solution}
